% ── Bring's Hierarchy ─────────────────────────────────────────────────────────

\chapter{Bring's Hierarchy}
\label{chap:brings-hierarchy}
\label{chap:brings-surface}% alias for backward compatibility

The \emph{Bring hierarchy} is the tower of algebraic curves
\[
  B_0 \;\hookrightarrow\; B_1 \;\hookrightarrow\; B_2 \;\hookrightarrow\; B_3,
\]
ordered by increasing genus $0, 0, 2, 4$ and connected by degree-$2$
branched coverings.  $B_3$ is Bring's curve proper; $B_2$ is its
hyperelliptic quotient; $B_1 = \mathbb{P}^1$ is the spacetime base;
and $B_0$ is the underlying ground field (a point in the Descartes
descent).  The hierarchy encodes a filtration of the QOU by topological
complexity and provides the geometric scaffolding for the particle
ansatz of Chapter~\ref{chap:path-integrals-braiding}.

% ── Section 1: General Definition ────────────────────────────────────────────

\section{General Definition and the Hierarchy}
\label{sec:hierarchy-general}

\begin{definition}[Bring hierarchy]
\label{def:bring-hierarchy}
The \emph{Bring hierarchy} is the sequence of smooth projective curves
\[
  B_0,\quad B_1,\quad B_2,\quad B_3
\]
over $\mathbb{C}$ (or any algebraically closed field $k$ of
characteristic $\neq 5$) with genera $g(B_n) = 2(n-1)$ for $n\ge 1$
and $B_0$ a point, together with degree-$2$ holomorphic maps
\[
  B_3 \xrightarrow{\;\pi_{32}\;} B_2
  \xrightarrow{\;\pi_{21}\;} B_1
  \xrightarrow{\;\pi_{10}\;} B_0,
\]
characterised as follows.  The covering $\pi_{21}\colon B_2\to B_1$ is
the hyperelliptic involution of $B_2$; the covering
$\pi_{32}\colon B_3\to B_2$ is the restriction of the degree-$4$ map
$\pi\colon B_3\to B_1$ to the intermediate quotient $B_2$.
\end{definition}

\begin{remark}[$B_0$: the point]
\label{rem:B0-point}
$B_0 = \mathrm{Spec}(k)$ is a single $k$-rational point.  In the QOU
descent picture (Remark~\ref{rem:descent-to-zero}), passing to $B_0$
means projecting out all spatial and chiral degrees of freedom; the
resulting object retains only the substrate parameter $q \in [1,\infty)$
as data.  Topologically, $B_0$ has genus $g = -1$ if one follows the
recursion $g(B_n)=2(n-1)$; we set $g(B_0)=\varnothing$ (undefined) and
regard $B_0$ as the degenerate base of the tower.
\end{remark}

\begin{remark}[$B_1 = \mathbb{P}^1(k)$, genus $0$]
\label{rem:B1-projective-line}
$B_1 = \mathbb{P}^1(k)$ is the projective line over $k$, the unique
smooth projective curve of genus $0$ with a $k$-rational point.  In
the QOU framework it plays the role of the \emph{spacetime base}: the
coordinate $z = x_0/x_1 \in \mathbb{P}^1$ parametrises the
one-dimensional spacetime direction inherited from the Descartes
descent.  The Euler characteristic is $\chi(B_1) = 2$.

The map $\pi_{10}\colon B_1 \to B_0$ is the structure morphism; it
records the fact that $\mathbb{P}^1$ has a unique rational point over the
ground field, namely the class $[B_1]$ itself.
\end{remark}

\begin{remark}[$B_2$, genus $2$]
\label{rem:B2-genus2}
$B_2$ is the unique (up to isomorphism over the moduli space) smooth
projective curve of genus $2$ that arises as the \emph{discriminant cover}
of $B_1$ branched at the six singular fibres of the elliptic fibration
of $B_3$ over $B_1$.  Its Euler characteristic is
$\chi(B_2) = 2 - 2\cdot 2 = -2$.

Concretely, if $\Delta_{B_3}(z)$ denotes the discriminant of the chiral
elliptic cubic of $B_3$
(equation~\eqref{eq:weierstrass-depressed}, Remark~\ref{rem:t-duality-elliptic}),
then $B_2$ is the hyperelliptic curve
\begin{equation}\label{eq:B2-hyperelliptic}
  B_2\colon\quad w^2 = \Delta_{B_3}(z),
\end{equation}
where $\Delta_{B_3}(z) = -4\,p(z)^3 - 27\,q(z)^2$, with
$p(z) = \tfrac{1}{3}(2z^2+z+2)$ and
$q(z) = \tfrac{5}{27}(4z^3+3z^2+3z+4)$, is a degree-$6$ polynomial
with $6$ simple roots (the branch locus of $\pi_{21}$).
The automorphism group of a generic $B_2$ is $\mathbb{Z}/2\mathbb{Z}$
(the hyperelliptic involution $\iota_2\colon(z,w)\mapsto(z,-w)$);
the $B_2$ arising here inherits additional automorphisms from the
$S_5$-symmetry of $B_3$.
\end{remark}

\begin{remark}[$B_3$: Bring's surface, genus $4$]
\label{rem:B3-brings-surface}
$B_3$ is \emph{Bring's curve}, the unique smooth projective curve of
genus $g = 4$ and Euler characteristic $\chi = -6$ defined by the
simultaneous vanishing of the first three power sums on $\mathbb{P}^4$
(Definition~\ref{def:brings-surface}, Section~\ref{sec:brings-surface}).
Its automorphism group $\mathrm{Aut}(B_3) \cong S_5$ has order $120$,
the maximum possible for a genus-$4$ curve by the Hurwitz bound.

The degree-$4$ map $\pi\colon B_3 \to B_1$ is a branched covering with
$6$ branch points at the roots of $\Delta_{B_3}(z)$; it factors as
$\pi = \pi_{21} \circ \pi_{32}$, where $\pi_{32}$ is a degree-$2$ map
to the intermediate curve $B_2$.
\end{remark}

\begin{proposition}[Riemann--Hurwitz for the hierarchy]
\label{prop:rh-hierarchy}
The three branched covers in the Bring hierarchy satisfy:
\begin{alignat}{3}
  &\pi_{32}\colon B_3\to B_2,\quad\deg=2,
  &&\quad 2\cdot 4 - 2 = 2(2\cdot 2-2)+2,
  &&\quad\text{$2$ simple ramification points,}
  \label{eq:rh-32}\\[4pt]
  &\pi_{21}\colon B_2\to B_1,\quad\deg=2,
  &&\quad 2\cdot 2 - 2 = 2(2\cdot 0-2)+6,
  &&\quad\text{$6$ simple branch points,}
  \label{eq:rh-21}\\[4pt]
  &\pi_{10}\colon B_1\to B_0,\quad\deg=1,
  &&\quad \text{unramified.}
  &&
  \label{eq:rh-10}
\end{alignat}
\end{proposition}
\begin{proof}
Equation~\eqref{eq:rh-21}: $\chi(B_2) = 2\chi(B_1) - R$ with $R=6$
ramification points gives $-2 = 4 - 6$ \checkmark.
Equation~\eqref{eq:rh-32}: $\chi(B_3) = 2\chi(B_2) - R$ with $R=2$
gives $-6 = -4 - 2$ \checkmark.  Equation~\eqref{eq:rh-10} is immediate.
\end{proof}

\begin{remark}[Genera and Euler characteristics]
\label{rem:hierarchy-genera}
The Bring hierarchy levels have the following topological invariants:
\begin{center}
\renewcommand{\arraystretch}{1.3}
\begin{tabular}{ccccc}
\hline
Level & Curve & Genus $g$ & $\chi = 2-2g$ & Automorphisms \\
\hline
$B_0$ & point & --- & --- & trivial \\
$B_1$ & $\mathbb{P}^1(k)$ & $0$ & $2$ & $\mathrm{PGL}(2,k)$ \\
$B_2$ & hyperelliptic & $2$ & $-2$ & $\mathbb{Z}/2\mathbb{Z}$ (generic) \\
$B_3$ & Bring's curve & $4$ & $-6$ & $S_5$ (order $120$) \\
\hline
\end{tabular}
\end{center}
\end{remark}

% ── Section 2: Coordinate Development for B_2 ────────────────────────────────

\section{Coordinate Development for $B_2$}
\label{sec:B2-coordinates}

We develop the coordinate structure of $B_2$ in parallel with the
treatment of $B_3$ in Section~\ref{sec:brings-surface}.  The key
simplification is that $B_2$ is a hyperelliptic curve: all its
geometry is determined by the pair $(z, w)$ satisfying
equation~\eqref{eq:B2-hyperelliptic}, plus the involution
$\iota_2\colon w \mapsto -w$.

\begin{definition}[$B_2$ as discriminant cover]
\label{def:B2-discriminant-cover}
Let $p(z) = \tfrac{1}{3}(2z^2+z+2)$ and
$q(z) = \tfrac{5}{27}(4z^3+3z^2+3z+4)$ be the
Weierstrass coefficients of the $B_3$ chiral elliptic fibration
(equation~\eqref{eq:weierstrass-depressed}).  The
\emph{discriminant curve} $B_2$ is the smooth projective model of the
affine curve
\begin{equation}\label{eq:B2-affine}
  w^2 \;=\; \Delta(z)
  \;:=\; -4\,p(z)^3 - 27\,q(z)^2
  \;\in\; \mathbb{Z}[z],
\end{equation}
a degree-$6$ polynomial.  The six simple roots $z_1,\dots,z_6$
of $\Delta(z)$ are the \emph{branch points} of $\pi_{21}\colon B_2\to B_1$
and coincide with the six singular fibres of the $B_3$ elliptic
fibration.
\end{definition}

\begin{remark}[QOU $4$-tuple for $B_2$]
\label{rem:B2-qou-tuple}
The $\mathbb{C}$ quantum observable universe
(Example~\ref{def:c-quantum-observable-universe}) for $B_2$
is the $4$-tuple $(M_{B_2}, G, P_{B_2}, E_{B_2})$ constructed as follows.

\medskip
\noindent\textbf{1.~CR manifold $M_{B_2}$.}\;
The unit circle bundle $M_{B_2} = \{(p,\eta)\in T^*B_2 : |\eta|_{g}=1\}$
over $B_2$ is a $3$-dimensional strictly pseudoconvex CR manifold
with contact form $\theta_{B_2} = \eta\,dz|_{M_{B_2}}$.

\medskip
\noindent\textbf{2.~Principal bundle $P_{B_2}$.}\;
The degree-$2$ covering $\pi_{21}\colon B_2\to\mathbb{P}^1$ branched
at the six roots $z_1,\dots,z_6$ lifts to a principal
$\mathrm{SU}(2)$-bundle $P_{B_2}\to M_{B_2}$.  The transition functions
are the $\mathrm{SU}(2)$ holonomy matrices around the six branch cuts;
the deck transformation $(z,w)\mapsto(z,-w)$ acts on $P_{B_2}$ by the
central element $-I\in\mathrm{SU}(2)$.

\medskip
\noindent\textbf{3.~State bundle $E_{B_2}$.}\;
The associated state bundle $E_{B_2} = P_{B_2}\times_G\mathbb{CP}^1$
has fibre $\mathbb{CP}^1$ over each point of $M_{B_2}$.  The two sheets
of $\pi_{21}$ give the two hemispheres of each $\mathbb{CP}^1$ fibre,
glued by the monodromy involution.

\medskip
\noindent\textbf{4.~Substrate parameter.}\;
The substrate parameter
$q(p) = 1 + |\langle L_{\theta_{B_2}}, F_{A}^+\rangle|(p)$
is maximal at the six branch points $z_i$ and tends to $1$ away from
them; the deep-quantum regime is therefore localised at
$\Delta(z)=0$.
\end{remark}

\begin{remark}[Holomorphic differentials on $B_2$]
\label{rem:B2-holomorphic-differentials}
A smooth genus-$2$ curve has a $2$-dimensional space of holomorphic
differentials $H^0(B_2, \Omega^1_{B_2})$.  For the hyperelliptic model
$w^2 = \Delta(z)$ the two basis differentials are
\begin{equation}\label{eq:B2-differentials}
  \omega_1 \;=\; \frac{dz}{w}
  \;=\; \frac{dz}{\sqrt{\Delta(z)}},
  \qquad
  \omega_2 \;=\; \frac{z\,dz}{w}
  \;=\; \frac{z\,dz}{\sqrt{\Delta(z)}}.
\end{equation}
These replace the single holomorphic differential of an elliptic curve;
the extra factor of $z$ in $\omega_2$ reflects the additional handle
(genus $2$ vs.\ genus $1$).  The period matrix
$\tau = (\oint_{A_j}\omega_i)^{-1}(\oint_{B_j}\omega_i)$ encodes the
complex structure of $B_2$; its imaginary part is positive definite
(Riemann's bilinear relations).
\end{remark}

\begin{remark}[Deck transformation and $\mathbb{Z}/2\mathbb{Z}$ spectral action]
\label{rem:B2-deck-transformation}
The hyperelliptic involution
$\iota_2\colon (z,w) \mapsto (z,-w)$
is the deck transformation of $\pi_{21}\colon B_2\to B_1$.  It acts on
the holomorphic differentials by
\[
  \iota_2^*\omega_1 = \omega_1,
  \qquad
  \iota_2^*\omega_2 = \omega_2,
\]
preserving the entire space $H^0(B_2,\Omega^1)$ (which is expected:
the deck transformation of the hyperelliptic map always preserves
the canonical bundle).  In terms of the DFT analogy with $B_3$:
whereas $B_3$ carries a $C_5$ spectral action
$w_k\mapsto\zeta_5^{-k}w_k$ on its Lagrange resolvents, $B_2$ carries a
$C_2 = \mathbb{Z}/2\mathbb{Z}$ action $w\mapsto -w$ corresponding to
the fourth root of unity $\zeta_4^2 = -1$.  The role of the DFT basis
$\{v_k\}_{k=0}^{3}$ with $v_k = \tfrac{1}{2}\sum_{j=0}^{3}i^{kj}y_j$
is played by the pair $(z, w)$ after the covering identification.
\end{remark}

\begin{remark}[Spectral forms for $B_2$]
\label{rem:B2-spectral-forms}
Parallel to the Lagrange resolvent spectral forms of $B_3$
(Remark~\ref{rem:lagrange-resolvents}), the $B_2$ model carries two
spectral constraints.  Setting $v_0 = \tfrac{1}{2}(y_0+y_1+y_2+y_3)$,
$v_1 = \tfrac{1}{2}(y_0+iy_1-y_2-iy_3)$,
$v_2 = \tfrac{1}{2}(y_0-y_1+y_2-y_3)$,
$v_3 = \tfrac{1}{2}(y_0-iy_1-y_2+iy_3)$ for homogeneous coordinates
$y_0,\dots,y_3$ on $\mathbb{P}^3$, the power-sum constraints
$\sum y_j = 0$ and $\sum y_j^2 = 0$ become:
\begin{align}
  v_0 &= 0
    &&\text{(linear projection),}
    \label{eq:B2-resolvent-linear}\\[4pt]
  2v_1 v_3 + v_2^2 &= 0
    &&\text{(quadric constraint).}
    \label{eq:B2-resolvent-quadric}
\end{align}
Together these cut out a smooth conic in $\mathbb{P}^2$,
isomorphic to $\mathbb{P}^1$, which is the spacetime base $B_1$.  The
full $B_2$ arises by adjoining the square root $w = \sqrt{\Delta(z)}$,
passing to the branched double cover of this $\mathbb{P}^1$.
The $C_4$ cyclic action $y_j\mapsto y_{j+1\bmod 4}$ descends to
the phase rotation $v_k\mapsto i^{-k}v_k$ in the spectral basis,
paralleling the $C_5$ rotation $w_k\mapsto\zeta_5^{-k}w_k$ of $B_3$.
\end{remark}

\begin{remark}[Color tube coordinates for $B_2$]
\label{rem:B2-color-tube}
The chiral fiber of $B_3$ has three roots $(x_2,x_3,x_4)$ giving three
color tube coordinates $(r,g,b)$
(Definition~\ref{def:color-tube-coordinates}).  For $B_2$ the analog is
a \emph{two-component} chiral fiber: at each $z\in B_1$, the two
branches of $\pi_{21}$ furnish two color coordinates
\begin{equation}\label{eq:B2-color-coords}
  \mathbf{r}^{(2)}_z = z + \sqrt{\Delta(z)},
  \qquad
  \mathbf{g}^{(2)}_z = z - \sqrt{\Delta(z)}.
\end{equation}
Their sum and product are
\begin{align}
  \mathbf{r}^{(2)}_z + \mathbf{g}^{(2)}_z &= 2z,
    \label{eq:B2-vieta-sum}\\[4pt]
  \mathbf{r}^{(2)}_z \cdot \mathbf{g}^{(2)}_z &= z^2 - \Delta(z).
    \label{eq:B2-vieta-product}
\end{align}
The $\mathbb{Z}/2\mathbb{Z}$ deck transformation swaps the two branches:
$\iota_2\colon (\mathbf{r}^{(2)}_z, \mathbf{g}^{(2)}_z) \mapsto (\mathbf{g}^{(2)}_z, \mathbf{r}^{(2)}_z)$.
This is the genus-$2$ analog of the $S_3$ permutation of three color
charges in $B_3$; the reduction from $S_3$ to $\mathbb{Z}/2\mathbb{Z}$
reflects the passage from three quark colors to two.
\end{remark}

\begin{proposition}[Chiral fiber of $B_2$ as $\mathbb{P}^1$]
\label{prop:B2-chiral-fiber}
For each $z \in B_1\setminus\{z_1,\dots,z_6\}$, the fiber
$(\pi_{21})^{-1}(z)$ consists of exactly two points, making the chiral
fiber of $B_2$ a genus-$0$ curve ($\mathbb{P}^1$).  At the six branch
points $z = z_i$ the fiber degenerates to a single (ramified) point.
\end{proposition}
\begin{proof}
Immediate from the definition of $\pi_{21}$ as the hyperelliptic cover
$w^2=\Delta(z)$: for $\Delta(z)\neq 0$ the equation has two solutions
$w=\pm\sqrt{\Delta(z)}$, while for $\Delta(z_i)=0$ the unique solution
is $w=0$.
\end{proof}

\begin{remark}[Braid group action on $B_2$]
\label{rem:B2-braid-action}
The mapping class group of the base $B_1 = \mathbb{P}^1$ punctured at
the six branch points $z_1,\dots,z_6$ is the braid group $B_6$ on six
strands.  Its action on $\pi_1$ of the punctured base lifts to the
\emph{symplectic representation}
\[
  \rho\colon B_6 \;\longrightarrow\; \mathrm{Sp}(4,\mathbb{Z}),
\]
which is the standard representation of the mapping class group
$\mathrm{Mod}_{0,6}$ on the homology $H_1(B_2,\mathbb{Z})\cong\mathbb{Z}^4$.
The monodromy around each branch point $z_i$ is a Dehn twist on $B_2$,
representing a simple transposition in $B_6$.  The full monodromy group
$\Gamma_2 = \mathrm{Sp}(4,\mathbb{Z})$ is the Siegel modular group, whose
orbifold quotient parametrises the principally polarised abelian surfaces
$J(B_2)$ (the Jacobians of $B_2$).  This replaces the $S_5$ automorphism
group of $B_3$ by a lattice group, reflecting the transition from finite
to infinite monodromy as we descend the hierarchy.
\end{remark}

\begin{remark}[Seiberg--Witten differential for $B_2$]
\label{rem:B2-SW-differential}
Parallel to the generation-dependent Seiberg--Witten differentials for
$B_3$ (Remark~\ref{rem:seiberg-witten}), the relevant
meromorphic $1$-form on $B_2$ is the \emph{$B_2$ Seiberg--Witten
differential}
\begin{equation}\label{eq:B2-SW-diff}
  \lambda_{B_2} \;=\; z\,\omega_1 \;=\; \frac{z\,dz}{w}
  \;=\; \frac{z\,dz}{\sqrt{\Delta(z)}}.
\end{equation}
Its $A$-cycle period
$a_{B_2} = \oint_{A_1}\lambda_{B_2}$
and dual $B$-cycle period
$a_{B_2}^D = \oint_{B_1}\lambda_{B_2}$
determine the effective prepotential of the genus-$2$ theory.
Because $B_2$ has genus $2$ rather than $1$, the prepotential carries
two independent moduli $(a_{B_2,1}, a_{B_2,2})$ corresponding to the
two handles of $B_2$; these are identified with the two lightest
generations in the QOU mass hierarchy.
\end{remark}

% ── Section 3: Bring's Surface ────────────────────────────────────────────────

% ── Bring's surface ──────────────────────────────────────────────────────────

\chapter{Bring's Surface}
\label{chap:brings-surface}

\begin{remark}[Descent to dimension zero]
\label{rem:descent-to-zero}
By the descent theorem (Theorem~\ref{thm:descartes}), it is sufficient
to study dimension-zero manifolds equipped with their bundles $P$ and~$E$.
At a point $p\in M_0$, the substrate parameter
$q(p)= 1 + \lvert\langle L_\theta,\, F_A^{+}\rangle\rvert(p)$
(Definition~\ref{def:substrate-parameter}) assigns a specific quantum
constant $q\in[1,\infty)$ determined by the principal bundle~$P$ and the
state bundle~$E$.  Here $E$ represents the totality of fermionic matter
packed into spacetime at the Big~Bang.  The trivial bundle $E=E_0$
corresponds to empty space, in which every point's location is exactly
known; a non-trivial bundle~$E$ begins to fill spacetime with fermions,
entering the regime of the quantum Bring surface whose underlying geometry
is non-commutative (Section~\ref{sec:quantum-brings}).
\end{remark}

\begin{definition}[Bring's surface]
\label{def:brings-surface}
Let $R$ be a commutative ring with $30 \in R^{\times}$ (i.e.\
$\mathrm{char}(R) \nmid 30$).  Let $x_j$, $j\in\{0,\dots,4\}$, be
homogeneous coordinates on $\mathbb{P}^4_R$.  \emph{Bring's surface} is the
algebraic variety
\[
  \Sigma_4 \;=\; \mathrm{Proj}\!\left(
    R[x_0,\dots,x_4]\Big/\bigl(\textstyle\sum x_j,\;
    \textstyle\sum x_j^2,\; \textstyle\sum x_j^3\bigr)
  \right),
\]
defined by the vanishing of the first three power sums.  Over any
algebraically closed field $K \supseteq R$, $\Sigma_4$ is a smooth
projective curve of genus~$4$ with automorphism group
$\mathrm{Aut}(\Sigma_4) \cong S_5$ ($\lvert\mathrm{Aut}\rvert = 120$,
achieving the Hurwitz bound $84(g-1) = 252$ up to the maximal value for
genus~$4$) and Euler characteristic
$\chi(\Sigma_4) = 2 - 2g = -6$.
\end{definition}

\begin{remark}[Realisations of Bring's surface]
\label{rem:brings-realisations}
Over an archimedean field $K$ (e.g.\ $\mathbb{Q}$, $\mathbb{R}$, or
$\mathbb{C}$), Bring's surface is a compact Riemann surface of genus~$4$.
The unit circle bundle $M_S$ in $T^*\Sigma_4$ is a CR manifold of
hypersurface type (real dimension~$3$), which provides the archimedean
realisation of the quantum observable universe
(Example~\ref{def:c-quantum-universe}).  Physics interpretations --- lepton
generations, chirality, colour --- are further realisations of the
$S_5$-symmetric structure (Remark~\ref{rem:descartes-leptons}).
\end{remark}

\begin{proposition}[Bring's surface is a quantum observable universe]
\label{prop:brings-is-qou}
Bring's surface $\Sigma_4$ (Definition~\ref{def:brings-surface}) gives rise
to a quantum observable universe over~$R$ (with $30 \in R^{\times}$):
\begin{enumerate}
  \item \emph{Quantum universe.}  Set $\mathcal{C} =
        \mathrm{Rep}(\mathrm{SU}_{q}(2))$ over~$R$,
        $A = \mathcal{O}_q(\mathrm{SU}(2))$ with Frobenius form given by the
        Haar state, and $\omega$ the forgetful functor to $R$-$\mathrm{Mod}$.
        The differential calculus $\Omega^1(A)$ is the Woronowicz
        $3$-dimensional calculus.
  \item \emph{Observability.}
        $\mathrm{Hom}(\mathbf{1},\, \mathrm{Ad}(A)) = 0$ by the quantum
        Schur lemma: $\mathrm{SU}_q(2)$ has no non-trivial invariants in the
        adjoint representation.
  \item \emph{Yekutieli lifting.}  The brane tower satisfies the Yekutieli
        lifting condition (Proposition~\ref{prop:yekutieli-lifting}).
  \item \emph{Euler characteristic.}  By Theorem~\ref{thm:qu-to-qou},
        $\chi(A) = \mathrm{tr}(\mathrm{id}_A) = 2 - 2g = -6$.
  \item \emph{Geometric descent} (for $n = 1$, $\dim M_S = 3$):
        \[
          M_S \;\xrightarrow{\;\mathrm{Reeb}\;}\; \Sigma_4
          \;\xrightarrow{\;\pi\;}\; \mathbb{P}(V)
          \;\xrightarrow{\;\mathrm{branch}\;}\; \{6\;\text{points}\}.
        \]
        The degree-$4$ covering $\pi$ has $6$ branch points with total
        ramification $R_{\mathrm{tot}} = 14$ (Riemann--Hurwitz), and the
        ramification profile is determined by the $S_5$-symmetry.
\end{enumerate}
\end{proposition}

\begin{proof}
\emph{Item~(1).}\;
The quantised coordinate ring $A = \mathcal{O}_q(\mathrm{SU}(2))$ is a
Hopf algebra over $R[q, q^{-1}]$ with antipode~$S$, counit~$\varepsilon$,
and the Haar state $h_q \colon A \to R$ furnishing a non-degenerate
Frobenius form $\lambda = h_q$ (non-degeneracy follows from the
faithfulness of the Haar state for compact quantum groups
\cite{woronowicz1987}).  The category
$\mathcal{C} = \mathrm{Rep}(\mathrm{SU}_q(2))$ is $R$-linear, rigid
monoidal (duals given by contragredient representations), and the forgetful
functor $\omega \colon \mathcal{C} \to R\text{-}\mathrm{Mod}$ is exact,
faithful, and monoidal.  The Woronowicz $3D$-calculus
$\Omega^1(A) = \ker(\varepsilon) / \ker(\varepsilon)^2$ is a
$3$-dimensional $A$-bimodule in~$\mathcal{C}$, providing the differential
calculus axiom.

\emph{Item~(2).}\;
The adjoint representation $\mathrm{Ad}(A)$ decomposes under
$\mathrm{SU}_q(2)$ into irreducible representations of spin $j \geq 1$
(the spin-$0$ component vanishes because
$\mathrm{Ad} \cong V_1 \oplus V_2 \oplus \cdots$ with no trivial
summand).  Therefore $\mathrm{Hom}(\mathbf{1}, \mathrm{Ad}(A)) = 0$ by
the quantum Schur lemma.

\emph{Item~(3).}\;
By Proposition~\ref{prop:jet-prolongation}, $J^1(A)$ is Frobenius--Hopf, so
the brane tower $\{A_n\}$ is well-defined.  Each transition
$A_n \hookrightarrow A_{n+1}$ gives a central extension of gerbes banded by
the centre $\mathcal{B} = Z(\mathrm{Aut}(A_n))$.  The obstruction class
$\kappa_n \in H^2(\mathcal{S}, \mathcal{B})$ exists by the standard gerbe
classification theorem, verifying the Yekutieli lifting condition
(Proposition~\ref{prop:yekutieli-lifting}).

\emph{Item~(4).}\;
The categorical trace of $\mathrm{id}_A$ in
$\mathcal{C} = \mathrm{Rep}(\mathrm{SU}_q(2))$ is the quantum dimension
$\mathrm{tr}(\mathrm{id}_A) = [N]_q \cdot [N-1]_q / [1]_q$ evaluated at
$N = 2$ (the rank of~$\mathrm{SU}(2)$).  At generic~$q$, this recovers
the topological Euler characteristic of~$\Sigma_4$:
$\chi(A) = 2 - 2g = 2 - 8 = -6$.

\emph{Item~(5).}\;
The Reeb quotient $M_S \to \Sigma_4$ collapses the $\mathrm{U}(1)$-orbits
of the circle bundle (weight-$0$ subalgebra $A_0 = A^{\mathrm{U}(1)}$).
The covering $\pi \colon \Sigma_4 \to \mathbb{P}(V)$ is the map to the
projective line induced by the linear system $|K_{\Sigma_4}|$; it has
degree~$4$ and genus~$4$ source, so by Riemann--Hurwitz
$R_{\mathrm{tot}} = 2g(\Sigma_4) - 2 - 4(2g(\mathbb{P}^1) - 2) =
6 - 4(-2) = 14$.
The $S_5$-symmetry of $\Sigma_4$ forces the $14$ units of ramification to
be distributed over exactly $6$ branch points (with ramification indices
determined by the cycle types of $S_5$).
\end{proof}

\begin{remark}[QOU $4$-tuple for Bring's surface]
\label{rem:brings-qou-4tuple}
We construct the archimedean quantum observable universe
(Example~\ref{def:c-quantum-observable-universe}) explicitly for Bring's
surface.  The five homogeneous coordinates $x_j$, $j\in\{0,\dots,4\}$,
split into \emph{spacetime} and \emph{chiral} degrees of freedom:
\[
  \underbrace{x_0,\, x_1}_{\text{spacetime}}
  \qquad
  \underbrace{x_2,\, x_3,\, x_4}_{\text{chiral}}.
\]
The projective ratio $z = x_0/x_1 \in \mathbb{CP}^1$ parametrises the
spacetime base.

\medskip
\noindent\textbf{1.\;CR manifold $M_S$.}\;
Let $S$ denote Bring's surface (Definition~\ref{def:brings-surface}).
Because $S$ is a Riemann surface (real dimension~$2$), it does not itself
carry a contact form.  The CR manifold is the unit circle bundle
$M_S = \{(p,\eta) \in T^*S : \lvert\eta\rvert_{g_S} = 1\}$,
where $g_S$ is the Fubini--Study metric restricted to~$S$.  The Liouville
$1$-form $\lambda = \eta\,dz$ on $T^*S$ restricts to a contact form
$\theta = \lambda\big|_{M_S}$, making $M_S$ a $3$-dimensional strictly
pseudoconvex CR manifold.

\medskip
\noindent\textbf{2.\;Principal bundle $P$.}\;
The branched covering $\pi\colon S \to \mathbb{CP}^1$ (degree $4$, six
branch points) lifts to a principal $\mathrm{SU}(2)$-bundle
$P \to M_S$.  In the spacetime coordinate~$z$, the transition functions
of~$P$ are the $\mathrm{SU}(2)$ holonomy matrices
$H_1,H_2$~\eqref{eq:holonomy-generators} around the branch cuts; in the
chiral coordinates $(x_2,x_3,x_4)$, the $S_3$ permutation of the three
sheets gives the $\mathrm{SO}(3)$ holonomy representation
(Definitions~\ref{def:su2-generators}--\ref{def:so3-generators}).

\medskip
\noindent\textbf{3.\;State bundle $E$.}\;
The associated state bundle is $E = P \times_G \mathbb{CP}^1$,
whose fibre at each point of $M_S$ is the chiral
$\mathbb{CP}^1$ parametrised by the three roots of the elliptic
cubic~\eqref{eq:chiral-elliptic}.

\medskip
\noindent\textbf{4.\;Connection, Levi form, and substrate parameter.}\;
The instanton connection $A$ on $P$ is expressed in the resolvent basis
$(w_1,\dots,w_4)$ (Remark~\ref{rem:lagrange-resolvents}) as the flat
$B_3$-monodromy connection whose holonomy around the six branch points is
generated by the trefoil holonomy matrices.  The Levi form on $M_S$ is
\[
  L \;=\; d\theta\big|_{H \times H},
  \qquad H = \ker\theta \subset TM_S,
\]
where $\theta = \eta\,dz\big|_{M_S}$ and the horizontal distribution $H$ is
spanned by the lifts of the chiral differentials $dx_2, dx_3, dx_4$ to
$M_S$.  The substrate parameter
(Definition~\ref{def:substrate-parameter}) then reads
\[
  q(p) \;=\; 1 + \lvert\langle L_\theta,\, F_A^+ \rangle\rvert(p),
  \qquad p \in M_S,
\]
coupling the chiral curvature to the contact geometry.  At the six branch
points (where chiral sheets merge), $\lVert F_A^+\rVert$ is maximal and
$q \to \infty$ (deep-quantum regime); away from the defects,
$q \to 1$ (classicality).

\medskip
\noindent\textbf{Summary.}\;
The $4$-tuple $(M_S,\, G,\, P,\, E)$, with $G = \mathrm{SU}(2)$,
constitutes an archimedean quantum observable universe
(Example~\ref{ex:brings-surface}), and every QOU structure---contact
form, Reeb field, Levi form, substrate parameter, $q$-deformed
operators---is determined by the spacetime coordinate~$z = x_0/x_1$ and
the chiral coordinates $(x_2, x_3, x_4)$.
\end{remark}

\begin{remark}[QOU axioms instantiated on Bring's surface]
\label{rem:brings-qou-axiom-map}
We record how each axiom of a quantum observable universe
(Definition~\ref{def:quantum-observable-universe}) is realised on Bring's
surface.

\medskip
\noindent\textbf{(i) Non-degeneracy.}\;
The Frobenius form $\lambda = h_q$ is the Haar state on
$A = \mathcal{O}_q(\mathrm{SU}(2))$.  Faithfulness of $h_q$ for compact
quantum groups~\cite{woronowicz1987} gives the isomorphism
$A \xrightarrow{\;\sim\;} A^{\vee}$.  The antipode $S$ and $h_q$
together define the dagger $f^{\dagger} = h_q \circ \mu(S(f) \otimes -)$,
which is positive-definite
(Definition~\ref{def:positive-definite-dagger}) because the Haar state
satisfies $h_q(a^* a) > 0$ for $a \neq 0$.

\medskip
\noindent\textbf{(ii) Observability.}\;
The adjoint representation decomposes as
$\mathrm{Ad}(A) \cong V_1 \oplus V_2 \oplus \cdots$ with no spin-$0$
summand, so $\mathrm{Hom}(\mathbf{1},\, \mathrm{Ad}(A)) = 0$ by the
quantum Schur lemma.  Geometrically, this is the statement that
Bring's surface has no global holomorphic vector field
(it has genus $g = 4 > 1$).

\medskip
\noindent\textbf{(iii) Categorical volume constraint.}\;
The fundamental object $V$ of
$\mathcal{C} = \mathrm{Rep}(\mathrm{SU}_q(2))$ has categorical
dimension~$2$.  In the archimedean realisation, $V$ is the rank-$2$
defining representation carried by the state bundle fibre
$E_p \cong \mathbb{C}^2$ (Remark~\ref{rem:brings-qou-4tuple}, item~3).
The non-degenerate epimorphism $\epsilon \colon \Lambda^2 V \to \mathbf{1}$
is instantiated by the Levi form
$L_\theta = d\theta|_{H \times H}$
(Remark~\ref{rem:volume-constraint-geometric}): because $L_\theta$ is
alternating and non-degenerate on the rank-$4$ real distribution
$H = \ker\theta$, it furnishes the symplectic pairing that makes $V$
self-dual and forces $\det(g_V) = 1$, selecting $\mathrm{SU}(2)$ over
$\mathrm{GL}(2)$.  In the spectral coordinates
(Remark~\ref{rem:lagrange-resolvents}), this constraint appears as the
Hermitian quadric $w_1 w_4 + w_2 w_3 = 0$, which is the $\mathrm{SL}(2)$
Pl\"ucker relation for $\Lambda^2 V$.  The top-degree extension to the
full chiral fibre is the spectral volume
constraint~\eqref{eq:spectral-volume-form}; see
Remark~\ref{rem:seiberg-witten} for the explicit symplectic
form~\eqref{eq:chiral-symplectic-form} and its relation to~$\epsilon$.
\end{remark}

\begin{remark}[Lagrange resolvents and spectral form]
\label{rem:lagrange-resolvents}
Let $\zeta=e^{2\pi i/5}$.  Apply the discrete Fourier transform to the
coordinates $x_j$ of Bring's surface
(Definition~\ref{def:brings-surface}):
\[
  w_k \;=\; \frac{1}{\sqrt{5}}\sum_{j=0}^{4}\zeta^{kj}\,x_j,
  \qquad k\in\{0,1,2,3,4\}.
\]
Under this change of basis the three defining power-sum equations become the
following \emph{spectral forms}:
\begin{align}
  w_0 &= 0
    &&\text{(linear projection),}
    \label{eq:resolvent-linear}\\[4pt]
  w_1\,w_4+w_2\,w_3 &= 0
    &&\text{(Hermitian quadric),}
    \label{eq:resolvent-quadric}\\[4pt]
  w_1^2\,w_3+w_4^2\,w_2+w_2^2\,w_1+w_3^2\,w_4 &= 0
    &&\text{(symmetric cubic).}
    \label{eq:resolvent-cubic}
\end{align}
The cyclic permutation $\sigma\colon x_j\mapsto x_{j+1}$ acts by
$\sigma(w_k)=\zeta^{-k}w_k$, so the $C_5$ symmetry is manifest as a phase
rotation $w_k\mapsto\zeta^k\,w_k$.
\end{remark}

\begin{remark}[Braid-group action on the resolvents]
\label{rem:braid-group-action}
The symmetric group $S_5$ acts on Bring's surface
(Remark~\ref{rem:brings-curve}) by permuting
$(x_0,\dots,x_4)$; this is the maximal automorphism group
(order~$120$) for genus~$4$.  Three distinguished families of $S_3$
subgroups govern the interplay between the base coordinates~$x_j$ and the
spectral resolvents~$w_k$
(Remark~\ref{rem:lagrange-resolvents}).  We describe each in terms of the
braid group (Definition~\ref{def:braid-group}), the resolvent coordinates,
and the $\mathrm{SU}(2)$/$\mathrm{SO}(3)$ matrices of
Definitions~\ref{def:su2-generators}--\ref{def:so3-generators}.

\medskip
\noindent\textbf{(i)~Spacetime action ($\mathbb{CP}^1$ base).}\;
Fix the chiral-like coordinates $x_2,x_3,x_4$ as internal fiber
degrees of freedom.  The constraint
$w_0=0$~\eqref{eq:resolvent-linear} forces $x_0+x_1=C$~(constant), so
the spacetime degrees of freedom collapse projectively to the ratio
$z=x_0/x_1\in\mathbb{CP}^1$.  The transposition
$x_0\leftrightarrow x_1$ acts as the M\"obius involution
$z\mapsto 1/z$, which is a $\pi$-rotation on $\mathbb{CP}^1$ and embeds
into $\mathrm{SU}(2)$ as $\beta=\exp(\pi\,e_1)$~\eqref{eq:su2-2cycle}.
In the resolvent basis this acts as
\begin{equation}\label{eq:spacetime-shift}
  w_k\;\longmapsto\;
  w_k+\tfrac{1}{\sqrt{5}}\,(1-\zeta^k)(x_1-x_0),
  \qquad k=1,\dots,4,
\end{equation}
an affine shift that fixes the chiral modes and translates the
$\mathbb{CP}^1$ fibre.

\medskip
\noindent\textbf{(ii)~Spectral $S_3$ (trefoil action).}\;
The frequency-coupled pairs $(w_1,w_4)$ and $(w_2,w_3)$ are permuted by
the twist and flip involutions
(Definition~\ref{def:braid-twist-flip}).
The full
$S_3$ structure emerges on the projective ratios
\[
  \eta_1=\frac{w_1}{w_2},\qquad\eta_2=\frac{w_3}{w_4}:
\]
$\sigma_1$ inverts both ratios $(\eta_1\mapsto\eta_1^{-1},\;
\eta_2\mapsto\eta_2^{-1})$, while $\sigma_2$ exchanges them
$(\eta_1\mapsto\eta_2^{-1},\;\eta_2\mapsto\eta_1^{-1})$.  The
$2$-cycle (reflection) swaps the pairs
$(w_1,w_4)\leftrightarrow(w_2,w_3)$; the $3$-cycle (rotation) sends
$(\eta_1,\eta_2)\mapsto(\eta_1,\,\zeta^j\eta_2)$ for $j\in\{1,2,3\}$,
permuting the three complex branches of the cubic
constraint~\eqref{eq:resolvent-cubic}.  The
$\mathrm{SU}(2)$ representation assigns the $3$-cycle to
$\alpha$~\eqref{eq:su2-3cycle} and the reflection to
$\beta$~\eqref{eq:su2-2cycle}; the $\mathrm{SO}(3)$
images~\eqref{eq:so3-generators} act as the holonomy on the spectral
energy distribution, with $R_\alpha$ rotating the time-like resolvent
pair~$(\omega_2,\omega_3)$ and $R_\beta$ reflecting the space-like
axis~$\omega_1$.

\medskip
\noindent\textbf{(iii)~Coordinate $S_3$.}\;
Fixing any two coordinates (say $x_0,x_1$) as a spacetime base and
permuting the remaining three $(x_2,x_3,x_4)$ gives a stabiliser
$\mathrm{Stab}(\{x_0,x_1\})\cong S_3\subset S_5$; there are
$\binom{5}{2}=10$ such subgroups.  In the resolvent basis a
permutation $\sigma\in S_3$ of $(x_2,x_3,x_4)$ acts as
\begin{equation}\label{eq:coord-s3-shift}
  w_k\;\longmapsto\;
  w_k+\tfrac{1}{\sqrt{5}}\sum_{j=2}^{4}
    \bigl(\zeta^{k\,\sigma(j)}-\zeta^{kj}\bigr)\,x_j,
\end{equation}
so the $w_k$ transform as vector fields over the $(x_0,x_1)$ plane and
the $S_3$ permutes how these fields couple to the base.  The same
$\mathrm{SO}(3)$ representation~\eqref{eq:so3-generators} realises this
action, now playing the r\^ole of the holonomy group on the internal
fibre.

\medskip
\noindent\textbf{Geometric $S_3$.}\;
When Bring's curve is viewed as a branched cover
$\mathcal{S}\to\mathbb{P}^1$, a third family of $S_3$ subgroups arises:
it permutes the three critical values of the covering map.  These three
families --- spectral, coordinate, and geometric --- correspond to the
three copies of $S_3$ inside $S_5$ identified in
Remark~\ref{rem:chirality-trefoil}.
% TODO: develop the monodromy representation as a quotient $B_3/S_3$
% and relate it to the branched-cover structure of the geometric $S_3$.
\end{remark}

\begin{definition}[Color tube coordinates]
\label{def:color-tube-coordinates}
  Let $z = x_0/x_1 \in \mathbb{CP}^1$ be the spacetime ratio
  (Remark~\ref{rem:braid-group-action}\textup{(i)}) and let $w_k$,
  $k \in \{0,\dots,4\}$, be the Lagrange resolvents of Bring's surface
  (Remark~\ref{rem:lagrange-resolvents}).  The \emph{color tube coordinates}
  are the shifted spectral coordinates
  \begin{equation}\label{eq:color-tube-coords}
    r = w_2 - z, \quad g = w_3 - z, \quad b = w_4 - z,
  \end{equation}
  together with the space-like shift $\sigma = w_1 - z$.  The
  \emph{color vector} is $\vec{c} = (r, g, b)$.
\end{definition}

\begin{remark}[Aggregate color quantities]
\label{rem:aggregate-color-quantities}
  Define the following aggregate quantities in color tube coordinates,
  with $\zeta = e^{2\pi i/5}$:
  \begin{align}
    \widetilde{S} &\;=\; \sigma + r + g + b
      &&\text{(total color deviation),}
      \label{eq:agg-S-tilde}\\[4pt]
    \widetilde{Q} &\;=\; \sigma b + r g
      &&\text{(paired color product),}
      \label{eq:agg-Q-tilde}\\[4pt]
    \widetilde{D} &\;=\; \zeta^{-1}\sigma + \zeta^{-2}r + \zeta^{-3}g + \zeta^{-4}b
      &&\text{(weighted color deviation),}
      \label{eq:agg-D-tilde}\\[4pt]
    \widetilde{C} &\;=\; \sigma^2 g + b^2 r + r^2\sigma + g^2 b
      &&\text{(intrinsic color cubic).}
      \label{eq:agg-C-tilde}
  \end{align}
  These are related to the unshifted aggregate quantities
  $S = \sum_{k=1}^{4}w_k$ and $D = \sum_{k=1}^{4}\zeta^{-k}w_k$ by
  \begin{equation}\label{eq:aggregate-shift}
    S \;=\; \widetilde{S} + 4z, \qquad D \;=\; \widetilde{D} - z,
  \end{equation}
  where the second identity uses $\sum_{k=1}^{4}\zeta^{-k} = -1$.
  The inverse Lagrange resolvent transform
  (Remark~\ref{rem:lagrange-resolvents}) gives
  \begin{equation}\label{eq:z-from-SD}
    z \;=\; \frac{x_0}{x_1} \;=\; \frac{S}{D}.
  \end{equation}
\end{remark}

\begin{remark}[Constraint equations in color tube coordinates]
\label{rem:color-tube-constraints}
  The Bring surface constraints~\eqref{eq:resolvent-linear}--\eqref{eq:resolvent-cubic}
  take the following form in color tube coordinates.

  \medskip
  \noindent\textbf{Spacetime recovery.}\;
  Substituting~\eqref{eq:aggregate-shift} into $z = S/D$
  (equation~\eqref{eq:z-from-SD}) yields $z(\widetilde{D} - z) = \widetilde{S} + 4z$,
  which rearranges to
  \begin{equation}\label{eq:spacetime-recovery}
    z^2 - (\widetilde{D} - 4)\,z + \widetilde{S} \;=\; 0.
  \end{equation}
  Thus the spacetime coordinate $z$ is determined implicitly by the
  shifted data $(\widetilde{S}, \widetilde{D})$.

  \medskip
  \noindent\textbf{Shifted quadric.}\;
  Substituting $w_k = c_k + z$ into the quadric~\eqref{eq:resolvent-quadric}
  gives
  \begin{equation}\label{eq:shifted-quadric}
    \widetilde{Q} + z\,\widetilde{S} + 2z^2 \;=\; 0,
  \end{equation}
  or equivalently $\sigma b + rg = -z(\sigma + r + g + b) - 2z^2$.

  \medskip
  \noindent\textbf{Shifted cubic.}\;
  Substituting $w_k = c_k + z$ into the cubic~\eqref{eq:resolvent-cubic}
  gives
  \begin{equation}\label{eq:shifted-cubic}
    \widetilde{C} + z(\widetilde{S}^2 - 2\widetilde{Q})
    + 3z^2\,\widetilde{S} + 4z^3 \;=\; 0.
  \end{equation}
  Eliminating $\widetilde{Q}$ via~\eqref{eq:shifted-quadric} yields the
  \emph{reduced cubic}:
  \begin{equation}\label{eq:reduced-cubic}
    \widetilde{C} + z\,\widetilde{S}^2 + 5z^2\,\widetilde{S} + 8z^3 \;=\; 0.
  \end{equation}
  All three equations are polynomial over $\mathbb{Z}$ in $z, r, g, b, \sigma$;
  they therefore define affine varieties over $\mathbb{Z}$ and descend
  to $\mathbb{F}_\pp$ for every prime $\pp \ne 5$.
\end{remark}

\begin{remark}[Interpretation of the color tube]
\label{rem:color-tube-interpretation}
  The terminology \emph{color tube} reflects two features.

  \medskip
  \noindent\textbf{Color.}\;
  The coordinates $r, g, b$ label the three chiral degrees of freedom
  that correspond, via the $S_3$ subgroups of $S_5$
  (Remark~\ref{rem:braid-group-action}\textup{(iii)}), to the three
  colour charges of quantum chromodynamics.  The chiral coordinates
  $x_2, x_3, x_4$ of the QOU $4$-tuple
  (Remark~\ref{rem:brings-qou-4tuple}) descend to $r = w_2 - z$,
  $g = w_3 - z$, $b = w_4 - z$ after shifting by the spacetime ratio~$z$.

  \medskip
  \noindent\textbf{Tube.}\;
  For fixed $z \in \mathbb{CP}^1$, the color vector $\vec{c} = (r, g, b)$
  is constrained by the shifted quadric~\eqref{eq:shifted-quadric} and
  cubic~\eqref{eq:shifted-cubic} to lie on a curve inside the
  three-dimensional color space $\mathbb{A}^3$.  As $z$ varies over the
  spacetime base, this curve sweeps out a tube --- the total space of
  the chiral fibration over $\mathbb{CP}^1$.

  \medskip
  \noindent\textbf{Space-like coordinate.}\;
  The shift $\sigma = w_1 - z$ is distinguished: it measures the deviation
  of the $C_5$-fundamental resolvent $w_1$ from the spacetime base,
  anchoring the $U(1)$ longitudinal backbone
  (Remark~\ref{rem:spectral-dynamics}) in the shifted frame.
\end{remark}

\begin{remark}[$C_5$ action in color tube coordinates]
\label{rem:c5-color-tube}
  The cyclic symmetry $w_k \mapsto \zeta^{-k} w_k$
  (Remark~\ref{rem:lagrange-resolvents}) becomes, in color tube coordinates:
  \begin{align}
    \sigma &\;\longmapsto\; \zeta^{-1}\sigma + (\zeta^{-1}-1)\,z,
      \label{eq:c5-sigma}\\[4pt]
    r &\;\longmapsto\; \zeta^{-2}r + (\zeta^{-2}-1)\,z,
      \label{eq:c5-r}\\[4pt]
    g &\;\longmapsto\; \zeta^{-3}g + (\zeta^{-3}-1)\,z,
      \label{eq:c5-g}\\[4pt]
    b &\;\longmapsto\; \zeta^{-4}b + (\zeta^{-4}-1)\,z.
      \label{eq:c5-b}
  \end{align}
  The $C_5$ action is no longer diagonal: each coordinate acquires an
  additive correction proportional to~$z$, reflecting the cost of
  breaking projective homogeneity by subtracting the affine
  quantity~$z$ from each resolvent.
\end{remark}

\begin{remark}[Intrinsic color cubic: pairing structure]
\label{rem:intrinsic-color-cubic-pairing}
  The intrinsic color cubic $\widetilde{C}$~\eqref{eq:agg-C-tilde}
  admits a natural decomposition into space-like and chiral sectors:
  \begin{equation}\label{eq:color-cubic-decomp}
    \widetilde{C}
    \;=\;
    \underbrace{\sigma(\sigma g + r^2)}_{\text{space--colour coupling}}
    \;+\;
    \underbrace{b(br + g^2)}_{\text{pure colour}}.
  \end{equation}
  The four terms pair the resolvent indices as $(1,3)$, $(2,1)$, $(4,2)$,
  $(3,4)$: each index is squared and contracted with the index two steps
  ahead in the cyclic order $1 \to 3 \to 4 \to 2 \to 1$.  This
  skip-one pairing respects the conjugation $w_k \leftrightarrow w_{5-k}$
  that pairs $(w_1, w_4)$ and $(w_2, w_3)$.
\end{remark}

\begin{proposition}[Chiral fibre as an elliptic curve]
\label{prop:chiral-elliptic}
Let $z=x_0/x_1\in\mathbb{CP}^1$ be the projective spacetime ratio
(Remark~\ref{rem:braid-group-action}\textup{(i)}).  Substituting the
constraint $x_0+x_1=C$ and eliminating $x_0,x_1$ in favour of~$z$,
the quadratic and cubic power-sum
conditions~\eqref{eq:resolvent-quadric}--\eqref{eq:resolvent-cubic} on
the chiral variable $X=(x_2,x_3,x_4)$ reduce at each fixed~$z$ to the
cubic
\begin{equation}\label{eq:chiral-elliptic}
  X^3+(z+1)\,X^2+(z^2+z+1)\,X+(z^3+z^2+z+1)=0.
\end{equation}
For generic~$z$ this defines an elliptic curve $E_z$ in the chiral
fibre; its $j$-invariant varies as a rational function of~$z$.  The
monodromy of~$E_z$ around the singular fibres (where the discriminant
of~\eqref{eq:chiral-elliptic} vanishes) recovers the holonomy
generators $H_1,H_2$~\eqref{eq:holonomy-generators} of
Remark~\ref{rem:holonomy-trefoil}.
\end{proposition}
\begin{proof}
Write $x_0=z\,x_1$ and use $x_0+x_1=C$ to eliminate both spacetime
variables.  The power-sum conditions $\sum x_j^k=0$ for $k=2,3$ then
become polynomial in~$z$ and $X=x_2+x_3+x_4$, $X_2=\sum_{i<j}x_ix_j$,
etc.  Reducing by the Newton identities yields the
cubic~\eqref{eq:chiral-elliptic}; its discriminant is a polynomial
in~$z$ whose roots give the singular fibres.
\end{proof}

\begin{remark}[Self-duality of the chiral elliptic curve]
\label{rem:t-duality-elliptic}
The cubic~\eqref{eq:chiral-elliptic} is symmetric under the exchange
$X\leftrightarrow z$: writing
\[
  F(X,z)\;=\;X^3+(z+1)\,X^2+(z^2+z+1)\,X+(z^3+z^2+z+1),
\]
one verifies $F(X,z)=F(z,X)$ term by term.  This involution
$\iota\colon(X,z)\mapsto(z,X)$ is a \emph{T-duality}: it identifies
every chiral fibre~$E_z$ with its dual spacetime wrapping~$E_X$, so
the internal and external degrees of freedom of the Bring surface are
algebraically indistinguishable.

\medskip\noindent\textbf{Quotient space.}\;
Passing to the symmetric coordinates $u=X+z$, $v=Xz$, the curve
becomes
\begin{equation}\label{eq:symmetric-quotient}
  v \;=\; \frac{(u+1)(u^2+1)}{2u+1}\,,
\end{equation}
which is rational in~$u$.  Hence the quotient
$E_z/\langle\iota\rangle\cong\mathbb{CP}^1$: folding the torus along
its $X=z$ diagonal collapses it to a genus-$0$ Riemann sphere.

\medskip\noindent\textbf{Weierstrass normal form.}\;
Depressing the cubic via $X=t-\tfrac{z+1}{3}$ eliminates the quadratic
term:
\begin{equation}\label{eq:weierstrass-depressed}
  t^{3}+p(z)\,t+q(z)=0,\qquad
  p=\frac{2z^2+z+2}{3},\quad
  q=\frac{5(4z^3+3z^2+3z+4)}{27}.
\end{equation}
The discriminant $\Delta(z)=-4\,p(z)^3-27\,q(z)^2$ is a polynomial in~$z$
whose roots determine the singular fibres of the elliptic
fibration $\mathcal{S}\to\mathbb{CP}^1$.  At those roots the
$j$-invariant $j(E_z)=-1728\cdot 64\,p(z)^3/\Delta(z)$ degenerates.
\end{remark}

\begin{remark}[Color tube cubic]
\label{rem:color-tube-cubic}
  At fixed $z \in \mathbb{CP}^1$, the three chiral roots of the
  elliptic cubic~\eqref{eq:chiral-elliptic} each satisfy the cubic in
  the color tube coordinate $\mathbf{c} \in \{\mathbf{r}_z, \mathbf{g}_z, \mathbf{b}_z\}$
  obtained by substituting $X = \mathbf{c} + z$
  into equation~\eqref{eq:chiral-elliptic}.  Expanding gives the
  \emph{color tube cubic}:
  \begin{equation}\label{eq:color-tube-cubic}
    \mathbf{c}^3 + (4z+1)\,\mathbf{c}^2 + (6z^2+3z+1)\,\mathbf{c}
    + (4z^3+3z^2+2z+1) \;=\; 0.
  \end{equation}
  This polynomial lies in $\mathbb{Z}[z][\mathbf{c}]$ and defines a
  variety over $\mathbb{Z}$.  The leading $z$-coefficients $4, 6, 4$
  equal $\binom{4}{1}$, $\binom{4}{2}$, $\binom{4}{3}$, and the
  constant term satisfies
  \[
    4z^3 + 3z^2 + 2z + 1 \;=\; \frac{d}{dz}(z^4 + z^3 + z^2 + z).
  \]

  \medskip
  \noindent Vieta's formulas for the three roots give:
  \begin{align}
    \mathbf{r}_z + \mathbf{g}_z + \mathbf{b}_z
      &\;=\; -(4z+1),
      \label{eq:vieta-sum}\\[4pt]
    \mathbf{r}_z\,\mathbf{g}_z + \mathbf{r}_z\,\mathbf{b}_z
      + \mathbf{g}_z\,\mathbf{b}_z
      &\;=\; 6z^2+3z+1,
      \label{eq:vieta-pairs}\\[4pt]
    \mathbf{r}_z\,\mathbf{g}_z\,\mathbf{b}_z
      &\;=\; -(4z^3+3z^2+2z+1).
      \label{eq:vieta-product}
  \end{align}
\end{remark}

\begin{definition}[Color tube fibers]
\label{def:color-tube-fibers}
  Let $z \in \mathbb{CP}^1$.  The \emph{color tube} at~$z$ is the
  affine curve defined over $\mathbb{Z}[z]$ by the color tube
  cubic~\eqref{eq:color-tube-cubic}; we denote it $R_z$, $G_z$, or
  $B_z$ according to which color tube coordinate $\mathbf{r}_z$,
  $\mathbf{g}_z$, or $\mathbf{b}_z$ is used:
  \begin{equation}\label{eq:color-tube-fiber-def}
    R_z \;:=\;
    \bigl\{\,\mathbf{r}_z \in \mathbb{A}^1 :
      \mathbf{r}_z^{\,3} + (4z+1)\,\mathbf{r}_z^{\,2}
      + (6z^2+3z+1)\,\mathbf{r}_z
      + (4z^3+3z^2+2z+1) = 0\,\bigr\},
  \end{equation}
  and $G_z$, $B_z$ are defined by the same equation with coordinates
  $\mathbf{g}_z$ and $\mathbf{b}_z$, respectively.  The coordinates
  are related to the Lagrange resolvents by
  \[
    \mathbf{r}_z = w_2 - z, \qquad \mathbf{g}_z = w_3 - z, \qquad
    \mathbf{b}_z = w_4 - z.
  \]
  As varieties, $R_z \cong G_z \cong B_z$ is isomorphic to the chiral
  elliptic fiber~$E_z$ of Proposition~\ref{prop:chiral-elliptic} via
  the shift $X = \mathbf{c} + z$.
\end{definition}

\begin{definition}[Observed spacetime value]
\label{def:observed-value}
  Let $\mathbf{r}_z, \mathbf{g}_z, \mathbf{b}_z$ be the color tube
  coordinates at a point of Bring's surface
  (Definition~\ref{def:color-tube-fibers}).  The
  \emph{observed spacetime value} is the element of $\mathbb{A}^1$
  recovered from the color tube coordinates by the Vieta
  identity~\eqref{eq:vieta-sum}:
  \begin{equation}\label{eq:observed-value}
    \mathbf{r}_z + \mathbf{g}_z + \mathbf{b}_z + 4z + 1 \;=\; 0.
  \end{equation}
  This identity defines the projection
  $\pi_z \colon R_z \times G_z \times B_z \to \mathbb{A}^1$,
  $({\mathbf{r}_z, \mathbf{g}_z, \mathbf{b}_z}) \mapsto z$,
  which is a polynomial map over $\mathbb{Z}$.  At a branch
  point $z_c$ where chiral sheets merge
  (Remark~\ref{rem:brings-smoothness}), the exceptional divisor is
  coordinatised by the limiting color vector; the observed spacetime
  value reads off $z_c$ from the color data via~\eqref{eq:observed-value},
  providing a coordinate on the exceptional $\mathbb{CP}^1$.
\end{definition}

\begin{remark}[T-duality in color tube coordinates]
\label{rem:t-duality-color-tube}
  The T-duality involution $\iota \colon (X, z) \mapsto (z, X)$
  (Remark~\ref{rem:t-duality-elliptic}) acts on color tube
  coordinates as follows.  Writing $c = X - z$, the dual point has
  base coordinate $z' = X = c + z$ and chiral coordinate $X' = z$,
  giving dual color tube coordinate $c' = z - X = -c$.  Hence
  \begin{equation}\label{eq:t-duality-color-tube}
    \iota \colon (c,\, z) \;\longmapsto\; (-c,\; c + z).
  \end{equation}
  The symmetry $F(X,z) = F(z,X)$
  (Remark~\ref{rem:t-duality-elliptic}) translates to covariance of
  the color tube cubic: the cubic~\eqref{eq:color-tube-cubic} at
  parameter~$z$ with root~$c$ is carried by $\iota$ to the cubic at
  parameter $c + z$ with root~$-c$.  One verifies this directly from
  $F(c+z,\, z) = F(z,\, c+z)$, since $F(z, c+z)$ is the color tube
  cubic at parameter $c+z$ evaluated at~$-c$.
\end{remark}

\begin{remark}[Sublocus $\widetilde{C}=0$ and colour confinement]
\label{rem:color-tube-subloci}
  \textbf{Vanishing of the intrinsic color cubic.}\;
  If $\widetilde{C} = 0$, the reduced cubic~\eqref{eq:reduced-cubic}
  factors as
  \begin{equation}\label{eq:reduced-cubic-factored}
    z\bigl(\widetilde{S}^2 + 5z\,\widetilde{S} + 8z^2\bigr) \;=\; 0,
  \end{equation}
  so either $z = 0$ (the spacetime origin) or $\widetilde{S}$ satisfies
  the quadratic $\widetilde{S}^2 + 5z\,\widetilde{S} + 8z^2 = 0$.
  The discriminant of this quadratic (in $\widetilde{S}$) is
  $25z^2 - 32z^2 = -7z^2$, which is negative for real $z \ne 0$.
  Over $\mathbb{C}$, the two solutions are
  \[
    \widetilde{S} \;=\; z\,\frac{-5 \pm \sqrt{-7}}{2};
  \]
  these are non-real for $z \ne 0$, so the locus $\widetilde{C} = 0$
  meets the real locus only at $z = 0$.  Over $\mathbb{F}_\pp$ the
  existence of a solution depends on whether $-7$ is a quadratic
  residue mod~$\pp$.

  \medskip
  \noindent\textbf{Pure colour sublocus.}\;
  The pure colour term $b(br + g^2)$
  (equation~\eqref{eq:color-cubic-decomp}) vanishes on
  \[
    \{b = 0\} \;\cup\; \{br + g^2 = 0\}.
  \]
  The condition $b = 0$ means $w_4 = z$: the fourth resolvent
  coincides with the spacetime ratio, collapsing one chiral degree of
  freedom onto the base.  The condition $br + g^2 = 0$, i.e.\
  $(w_4 - z)(w_2 - z) + (w_3 - z)^2 = 0$, is a quadratic relation
  among the chiral resolvents and~$z$.  Either condition can be
  interpreted as a \emph{colour-confinement} constraint: one chiral
  mode is frozen in register with the spacetime coordinate, and the
  remaining two are coupled by a Hermitian-quadric relation in the
  shifted frame.
\end{remark}

\begin{remark}[Chiral space as spectral curve and spinor bundle]
\label{rem:chiral-bundles}
The spacetime--chiral decomposition of Bring's surface maps naturally onto
three classical bundle structures over the $\mathbb{CP}^1$ base $z=x_0/x_1$.

\medskip
\noindent\textbf{Cotangent bundle.}\;
The three-sheeted cubic~\eqref{eq:chiral-elliptic} realises the chiral
variables $X=(x_2,x_3,x_4)$ as a \emph{spectral curve}
$\Sigma\subset T^{*}\mathbb{CP}^1$: the base coordinate~$z$ parametrises
spacetime position, and the three roots $X_1,X_2,X_3$ of
Proposition~\ref{prop:chiral-elliptic} are eigenvalues of a Higgs-like
field $\Phi$ on the base.  The six singular fibres (where the discriminant
$\Delta(z)$ vanishes) are the loci at which eigenvalues coincide and the
differential $dX$ degenerates in the cotangent fibre.

\medskip
\noindent\textbf{Spinor bundle.}\;
The spectral modes $w_1,\dots,w_4$
(Remark~\ref{rem:lagrange-resolvents}) live in a rank-$2$ complex vector
bundle over $\mathbb{CP}^1$.  The chiral doublet
$\Psi_L=(w_1,w_2)^T$ and its conjugate $\Psi_R=(w_4,w_3)^T$ transform
under the $\mathrm{SU}(2)$ holonomy generators
$H_1,H_2$~\eqref{eq:holonomy-generators} of
Remark~\ref{rem:holonomy-trefoil}, making $\Psi_L$ a section of the
\emph{spinor bundle} $\mathcal{S}^{+}\to\mathbb{CP}^1$.  The $B_3$
monodromy equips this bundle with a flat connection whose holonomy is the
trefoil-knot group.

\medskip
\noindent\textbf{Adjoint bundle.}\;
The $\mathrm{SO}(3)$ images $R_\alpha,R_\beta$~\eqref{eq:so3-generators}
act on the triple $(\omega_1,\omega_2,\omega_3)$, which therefore spans a
section of the adjoint bundle $\mathrm{ad}\,P\to\mathbb{CP}^1$ associated
to the principal $\mathrm{SU}(2)$-bundle~$P$.  The T-duality involution
$\iota\colon X\leftrightarrow z$ (Remark~\ref{rem:t-duality-elliptic})
interchanges the fibre and base of the adjoint bundle, exchanging
spectral eigenvalues with spacetime coordinates.
\end{remark}

\begin{remark}[Smoothness of the Bring surface and resolution of spacetime defects]
\label{rem:brings-smoothness}
The Bring surface $\mathcal{S}$ (Definition~\ref{def:brings-surface}) is a
\emph{smooth} algebraic variety: the three power-sum constraints are
transverse in~$\mathbb{P}^4$, so their intersection carries no singular
points.  The six topological defects identified in
Remark~\ref{rem:chiral-bundles}---where the discriminant $\Delta(z)$
vanishes and two chiral sheets $X_i,X_j$ collide---are singularities of
the \emph{projection} $\pi\colon\mathcal{S}\to\mathbb{CP}^1$, not of the
surface itself.

\medskip
\noindent\textbf{Resolution in the $x_j$ coordinates.}\;
At a branch point $z=z_c$ where, say, $x_2=x_3$, define the local chiral
coordinate $\delta=x_2-x_3$.  In a neighbourhood of the defect the
spacetime ratio satisfies
\[
  z-z_c\;\sim\;\delta^2,
\]
so while the projection map develops a ramification cusp, $\delta$ passes
smoothly through zero, carrying the geometry from one sheet to the other.
The exceptional divisor of the blow-up at each branch point is therefore
already embedded in the chiral coordinates $x_2,x_3,x_4$: no external
resolution is required.

\medskip
\noindent\textbf{Resolution in the $w_k$ coordinates.}\;
The DFT mixes all five $x_j$ into every resolvent $w_k$
(Remark~\ref{rem:lagrange-resolvents}), so no single $w_k$ vanishes at a
branch point.  The Briggs--Klein chart
(Definition~\ref{def:briggs-klein-chart}) remains a smooth coordinate
system throughout $0<r<\infty$: the radial parameter~$r=\lvert w_1\rvert^2$
is a regular function on $\mathcal{S}$, and the $C_5$ orbits
$w_k\mapsto\zeta^k w_k$ foliate the surface without fixed points.  The
infinite torsion and eigenvalue crossings visible in the base
$\mathbb{CP}^1$ are artefacts of projecting away the internal degrees of
freedom that the $w_k$ retain.
\end{remark}

\begin{definition}[Briggs--Klein affine chart]
\label{def:briggs-klein-chart}
Let $r=\lvert w_1\rvert^2$ denote the squared modulus of the first resolvent
coordinate on Bring's surface
(Definition~\ref{def:brings-surface}).  The
\emph{Briggs--Klein affine chart} is the coordinate system $(w_1,w_2)$
restricted to the open set $0<r<\infty$, equipped with the Hermitian metric
inherited from the Fubini--Study metric on $\mathbb{P}^4$.  In this chart
the $C_5$ action is the phase rotation $w_k\mapsto\zeta^k\,w_k$, and the
radial coordinate $r$ is $C_5$-invariant.
\end{definition}

\begin{remark}[Radial projection of Bring's surface]
\label{rem:radial-projection}
Restricting a Briggs--Klein affine chart
(Definition~\ref{def:briggs-klein-chart}) to the interval $r\in(0,1)$ and
projecting Bring's surface onto this radial coordinate yields a fibration
over the base manifold $M$ with trivial bundle $E=E_0$.  At each
$r$, the coordinate $w_1$ is restricted to a circle of radius $\sqrt{r}$ in
the complex plane $\mathbb{C}_{w_1}$.  This circle is the orbit of the $C_5$
symmetry: the rotation $w_1\mapsto\zeta\,w_1$ moves the point along this
circle without changing~$r$.
\end{remark}

\begin{remark}[Four-colourability of Bring's surface]
\label{rem:four-colourability}
By the Heawood bound, an orientable surface of genus~$g$ has chromatic number
at most $\lfloor(7+\sqrt{1+48g})/2\rfloor$.  For genus~$4$ this bound
equals~$10$.  Bring's surface, however, admits a proper colouring with only
four colours --- considerably fewer than the maximum permitted by its genus.
This low chromatic number reflects the high degree of symmetry: the $S_5$
automorphism group constrains the combinatorics of any cell decomposition,
forcing a more efficient colouring than a generic genus-$4$ surface would
require.
\end{remark}

\begin{remark}[The four-colour theorem and Bring's surface]
\label{rem:four-colour-theorem}
The classical four-colour theorem asserts that every planar graph (equivalently,
every graph embeddable in $S^2$) is four-colourable.  For higher-genus
surfaces the chromatic number can be much larger: the Heawood conjecture
(proved by Ringel--Youngs for $g\geq 1$) shows that the maximum chromatic
number grows as $O(\sqrt{g})$.  That Bring's surface attains the planar
bound of four --- despite living at genus~$4$ --- is a reflection of the
rigidity imposed by its $S_5$ symmetry.  In the language of the Descartes
universe, the trivial bundle $E=E_0$
(Remark~\ref{rem:radial-projection}) carries no torsion, and the resulting
flatness of the tessellation forces the same chromatic economy as the
sphere.
\end{remark}

\begin{remark}[Cyclotomic reduction and $C_{\mathfrak{p}}$ preservation]
\label{rem:cyclotomic-reduction}
Let $\mathfrak{p}$ be a prime (not to be confused with a point $p\in M$).
Working over the cyclotomic ring
\begin{equation}\label{eq:cyclotomic-ring}
  \mathbb{Z}_{\zeta_{\mathfrak{p}}}
  \;:=\;\mathbb{Z}[i,\,\zeta_{\mathfrak{p}}],
\end{equation}
where $\zeta_{\mathfrak{p}}=e^{2\pi i/\mathfrak{p}}$ is a primitive
$\mathfrak{p}$-th root of unity, consider the sublocus
\[
  S_{\mathfrak{p}}
  \;:=\;
  \Sigma_4\;\cap\;\mathbb{Z}_{\zeta_{\mathfrak{p}}}^{\,5},
\]
that is, the set of points on Bring's surface whose resolvent coordinates
lie in $\mathbb{Z}_{\zeta_{\mathfrak{p}}}$.  The cyclic subgroup
$C_{\mathfrak{p}}\subset S_5$ generated by the $\mathfrak{p}$-cycle preserves
$S_{\mathfrak{p}}$: for any $P\in S_{\mathfrak{p}}$, the image
$\sigma^k(P)$ remains in $\mathbb{Z}_{\zeta_{\mathfrak{p}}}^{\,5}$ because
$\sigma$ acts by multiplication by $\zeta_{\mathfrak{p}}^{-k}$ in the
resolvent basis (Remark~\ref{rem:lagrange-resolvents}), and
$\zeta_{\mathfrak{p}}^{-k}\in\mathbb{Z}_{\zeta_{\mathfrak{p}}}$.  This
arithmetic compatibility between the cyclic symmetry and the cyclotomic ring
underlies the descent of the $C_5$ action to finite characteristic.
\end{remark}

\begin{definition}[Quaternionic spectral resolvents]
\label{def:quaternionic-spectral-resolvents}
Let $\mathcal{S}$ denote the Bring's surface
(Definition~\ref{def:brings-surface}).  Let $q\in\mathbb{C}$ with
$0<\lvert q\rvert\le 1$, and set $r=-\ln\lvert q\rvert$ so that
$\lvert q\rvert=e^{-r}$.
Write $\mathbf{i},\mathbf{j},\mathbf{k}$ for the standard basis of
$\operatorname{Im}\mathbb{H}$.
Using the Lusztig $q$-exponential and $q$-logarithm,
\begin{equation}\label{eq:q-exp}
  \exp_q(x)\;=\;\sum_{n=0}^{\infty}\frac{x^n}{[n]_q!},
  \qquad
  \ln_q(1+x)\;=\;\sum_{n=1}^{\infty}(-1)^{n+1}\frac{x^n}{[n]_q},
\end{equation}
we define the \emph{quaternionic spectral resolvents}
$\omega_k\colon\mathcal{S}\to J^\infty E$ by
\begin{align}
  \omega_0 &\;=\;
    \Bigl(\tfrac{1}{\sqrt{5}}\!\sum_{j=0}^{4}x_j\Bigr)\,
    \ln_q(1+r)
    &&\text{(uncertainty),}
    \label{eq:omega-0}\\[4pt]
  \omega_1 &\;=\;
    \frac{\mathbf{i}}{\ln_q(1+r)}\,
    \Bigl(\tfrac{1}{\sqrt{5}}\!\sum_{j=0}^{4}\zeta^{j}\,x_j+i\Bigr)
    &&\text{(space-like),}
    \label{eq:omega-1}\\[4pt]
  \omega_2 &\;=\;
    \mathbf{j}\,
    \Bigl(\tfrac{1}{\sqrt{5}}\!\sum_{j=0}^{4}\zeta^{2j}\,x_j+i\Bigr)\,
    \bigl(\exp_q(r)-1\bigr)
    &&\text{(time-like),}
    \label{eq:omega-2}\\[4pt]
  \omega_3 &\;=\;
    \mathbf{k}\,
    \Bigl(\tfrac{1}{\sqrt{5}}\!\sum_{j=0}^{4}\zeta^{3j}\,x_j+i\Bigr)\,
    \bigl(\exp_q(r)-1\bigr)
    &&\text{(time-like).}
    \label{eq:omega-3}
\end{align}
By the first power sum~\eqref{eq:resolvent-linear}, $\omega_0=0$ in the
balanced state; the remaining coordinates embed $\mathcal{S}$ into the
infinite jet shell $J^\infty E$
(Remark~\ref{rem:jet-shell}).
\end{definition}

\begin{remark}[Algebraic constraints in the $\omega_k$ basis]
\label{rem:omega-constraints}
In the spectral resolvent basis
(Definition~\ref{def:quaternionic-spectral-resolvents}) the
Bring's surface equations take the following form:
\begin{itemize}
\item \emph{Barycentric:}\; $\omega_0=0$.
\item \emph{Quadric ($\mathbb{H}$ condition):}\;
  $\omega_1\bar{\omega}_1+\omega_2\bar{\omega}_3=0$.
  This preserves the horizontal self-duality $F_A|_H={*_H}F_A|_H$
  (Proposition~\ref{prop:self-dual-curvature}) and enforces a
  Minkowski-like balance between the space-like and time-like modes.
\item \emph{Cubic (phase locking):}\; a quintic resonance condition that
  locks the expansion of $\ln_q$ to $\exp_q$, preserving the $C_5$
  symmetry and the Descartes condition $\chi=-6$.
\end{itemize}
\end{remark}

\begin{remark}[Spectral dynamics and the $\mathrm{Sp}(1)$ action]
\label{rem:spectral-dynamics}
The growth profiles in
Definition~\ref{def:quaternionic-spectral-resolvents} separate the space-like
and time-like sectors of the Bring's surface:
\begin{itemize}
\item \emph{$\omega_1$ (space-like):}\; the factor $1/\ln_q(1+r)$ diverges
  as $r\to 0$, so $\omega_1$ achieves infinite density at the base stratum
  $M_0$, characterising the primordial vortex.  This mode anchors the
  $U(1)$ longitudinal backbone of the $\mathrm{SU}(2)$ bundle.
\item \emph{$\omega_2,\omega_3$ (time-like):}\; the factor
  $\exp_q(r)-1$ vanishes at $r=0$ and grows rapidly for small $r$,
  driven by the $q$-deformed Reeb flow
  (Definition~\ref{def:q-reeb-flow}).  These ladder modes expand at the big
  bang before tapering as the geometric viscosity
  (Remark~\ref{rem:q-harmonic-physical}) takes hold.
\end{itemize}
To realise the rotation as a function of $r$, embed the radial parameter
into $\mathrm{Sp}(1)$ via the unit quaternion
\begin{equation}\label{eq:sp1-quaternion}
  \mathbf{q}(r)
  \;=\;
  \cos_q\!\theta(r)+\mathbf{u}\,\sin_q\!\theta(r),
  \qquad
  \theta(r)=\tfrac{\pi}{2}\,e^{-r},
\end{equation}
where $\cos_q$ and $\sin_q$ are the $q$-deformed trigonometric functions
(built from $\exp_q$, equation~\eqref{eq:q-exp})
\[
  \cos_q(x)=\frac{\exp_q(ix)+\exp_q(-ix)}{2},
  \qquad
  \sin_q(x)=\frac{\exp_q(ix)-\exp_q(-ix)}{2i},
\]
and $\mathbf{u}\in\operatorname{Im}\mathbb{H}$ is a fixed unit imaginary
quaternion.  The conjugation
$\boldsymbol{\omega}'=\mathbf{q}(r)\,\boldsymbol{\omega}\,\mathbf{q}(r)^{*}$
gives maximum twist ($\theta=\pi/2$) in the deep-quantum regime $r\to 0$
and classical freeze ($\theta\to 0$) as $r\to\infty$.
\end{remark}

\begin{remark}[Primordial vortex and the deep-quantum limit]
\label{rem:primordial-vortex}
As $r\to 0$ the spectral resolvents converge toward the
\emph{primordial vortex}:
\begin{itemize}
\item The quantum Planck constant
  $\hbar_q\to 1$ (Definition~\ref{def:quantum-planck-constant}).
\item The $q$-torsion leaves
  (Definition~\ref{def:q-torsion-leaf}) collapse to the zero-dimensional
  stratum $M_0=\{\pm 1\}$.
\item The $\omega_0$ coordinate, although identically zero in the
  balanced state, formally blows up in the jet prolongation, representing
  the concentration of energy at a singular spatial point.
\end{itemize}
Space-like energy is trapped because the time-like rotation becomes infinite
in frequency but zero in displacement, effectively turning the
$\mathrm{SU}(2)$ holonomy into a localised topological charge.
\end{remark}

\begin{remark}[Geometric drag and the mass residue]
\label{rem:geometric-drag}
The Levi form $L$ (Definition~\ref{def:substrate-parameter}) must be
preserved under the $\mathrm{Sp}(1)$ conjugation.  Since
\[
  \lVert\boldsymbol{\omega}'\rVert^2
  =\lVert\boldsymbol{\omega}\rVert^2
  =\omega_1\bar{\omega}_1+\omega_2\bar{\omega}_2+\omega_3\bar{\omega}_3,
\]
the quadric constraint~\eqref{eq:resolvent-quadric} remains satisfied and
the horizontal self-duality $F_A|_H={*_H}F_A|_H$ is preserved.  The mass of a stable vortex is
recovered as the residue of the interaction between the growth profiles:
\begin{equation}\label{eq:mass-residue}
  m \;\propto\;
  \oint_{\gamma}
  \lVert\nabla_A q\rVert\cdot
  \lVert\Omega_1\wedge\Omega_{2,3}\rVert,
\end{equation}
where $\Omega_k$ denotes the curvature form restricted to the $\omega_k$
sector.  This heuristic is made rigorous by the nonabelian multiplicative
integral of the mass identity (Definition~\ref{def:mass-identity}).
When $q=1$ (trivial bundle), the $\omega_k$ commute and $m=0$;
when $q\to 0$ (maximal twist), energy is trapped in the $\omega_1$ spatial
filament.
\end{remark}

% ── Geometric resolution and the Seiberg-Witten differential ─────────────────

\begin{remark}[Geometric resolution and the Seiberg--Witten differential]
\label{rem:seiberg-witten}
The infinite geometric torsion~$T$ and the six topological defects on the
$\mathbb{CP}^1$ base $z=x_0/x_1$ are strictly artefacts of
dimensional projection.  The total space, when lifted into the
complete coordinate framework, is a globally smooth complex manifold.

\medskip
\noindent\textbf{1.~The blow-up and the exceptional divisor.}\;
Near a spacetime defect $z_c$ where chiral sheets (e.g.\ $x_2$ and
$x_3$) merge, the singularity is resolved by blowing up the point into
an exceptional divisor parameterised by the local coordinate
$\delta=x_2-x_3$.  The local geometry satisfies
\[
  z-z_c\;\sim\;\delta^2,
\]
proving that the spacetime base ``bounces'' at the defect: the chiral
coordinate continuously glides through the exceptional divisor and
transitions between sheets without encountering a singularity
(cf.\ Remark~\ref{rem:brings-smoothness}).

In the spectral domain, the Jacobian determinant of the DFT
(Remark~\ref{rem:lagrange-resolvents}) is a constant
$\lvert\det(J)\rvert=1$, so the global volume form never collapses:
\begin{equation}\label{eq:spectral-volume-form}
  \Omega \;=\; dw_1\wedge dw_2\wedge dw_3\wedge dw_4
  \;=\; \det(J)\;dx_0\wedge\cdots\wedge dx_4\Big|_{w_0=0}.
\end{equation}
This is the \emph{spectral volume constraint}: the top-degree analogue
of the categorical volume constraint
(Definition~\ref{def:categorical-volume-constraint}).  The categorical
constraint provides a non-degenerate pairing
$\epsilon \colon \Lambda^2 V \to \mathbf{1}$ on the rank-$2$ fundamental
object; the spectral volume constraint
extends this to the full chiral fibre, ensuring that the $4$-form~$\Omega$
is everywhere non-degenerate.  Both enforce $\det = 1$: the categorical
form forces $g_V \in \mathrm{SL}(V)$, while the spectral form forces
$\det(J) \in \mathrm{U}(1)$.

The link is made explicit by the symplectic form.  On the categorical
side, the alternating pairing on $V \cong \mathbb{C}^2$ is
(Lemma~\ref{lem:volume-from-rigidity})
\[
  \epsilon(v \wedge w)
  \;=\; v_1 w_2 - v_2 w_1,
  \qquad v, w \in V.
\]
On the Brings surface, the spectral coordinates split into conjugate
pairs $(w_1, w_4)$ and $(w_2, w_3)$ with
$\overline{w_k} = w_{5-k}$.  The quadric
constraint~\eqref{eq:resolvent-quadric} is
$w_1 w_4 + w_2 w_3 = 0$,
which is the Pl\"ucker relation $\Lambda^2 V \hookrightarrow
V \otimes V$.  The induced symplectic $2$-form on the chiral fibre is
\begin{equation}\label{eq:chiral-symplectic-form}
  \omega_{\mathrm{ch}}
  \;=\; dw_1 \wedge dw_4 + dw_2 \wedge dw_3,
\end{equation}
and $\Omega = \tfrac{1}{2}\,\omega_{\mathrm{ch}} \wedge
\omega_{\mathrm{ch}}$, so the non-degeneracy of $\Omega$
(i.e.\ $\det(J) \neq 0$) is equivalent to the non-degeneracy of
$\omega_{\mathrm{ch}}$, which is exactly the fibre-functor image of the
categorical volume form~$\epsilon$.

The geometric torsion~$T$ is thereby localised not as a point
singularity but as a trapped, finite magnetic flux integrated over
the surface of the exceptional $\mathbb{CP}^1$ sphere.

\medskip
\noindent\textbf{2.~The Seiberg--Witten differential.}\;
With the singularities geometrically resolved, the dynamics of the chiral
space are completely parameterised by a meromorphic $1$-form embedded in
the cotangent bundle $T^{*}\mathbb{CP}^1$, the \emph{Seiberg--Witten
differential}:
\begin{equation}\label{eq:seiberg-witten-differential}
  \lambda_{\mathrm{SW}} \;=\; X\,dz,
\end{equation}
where $X\in\{x_2,x_3,x_4\}$ is the internal momentum (Higgs-field
eigenvalue) defining the spectral curve
(Remark~\ref{rem:chiral-bundles}), and $z$ is the spacetime position.
The exact physical properties of the topological defects---specifically
the mass/energy spectrum generating the torsion---are rigorously
determined by integrating this differential along the fundamental
non-contractible loops ($A$ and $B$ cycles) of the resolved chiral
torus:
\begin{equation}\label{eq:seiberg-witten-periods}
  a \;=\; \oint_{A}\lambda_{\mathrm{SW}},
  \qquad
  a_D \;=\; \oint_{B}\lambda_{\mathrm{SW}}.
\end{equation}

By lifting the Bring surface into the full spatial and spectral
coordinate spaces, the degenerate branch points resolve completely into
smooth exceptional divisors.  The infinite local torsion is geometrically
captured as trapped flux, revealing that the discrete braid-group
($B_3$) holonomy (Remark~\ref{rem:holonomy-trefoil}) is simply the
continuous Kodaira monodromy of a globally smooth Seiberg--Witten gauge
theory.
\end{remark}



\begin{remark}[Reeb vector field on the exceptional divisor]
\label{rem:reeb-exceptional-divisor}
The Reeb vector field $\xi_q$
(Definition~\ref{def:q-reeb-vector-field}) is the dynamical
complement to the static topology of the Bring surface.  While the
surface itself provides the handles (genus~$4$), the Reeb vector
provides the ``wind'' blowing through them, and the quantum torsion~$q$
sets the strength of that wind.

\medskip
\noindent\textbf{The classical limit ($q\to 1$): cones and frozen geometry.}\;
If the substrate parameter $q\to 1$
(Definition~\ref{def:substrate-parameter}), the
torsion coupling vanishes and the phase rotation
$w_k\mapsto\zeta^k\,w_k$ loses its temporal component.  The fibre
over each base point ceases to twist as one moves along a path on
$\mathbb{CP}^1$: the chiral tubes are frozen in the configuration
determined by the power-sum equations.  Metrically this is a conical
limit---the curvature concentrates at the blow-up points but no
circulating torsion field smooths it out.  In terms of the Reeb
flow (Definition~\ref{def:q-reeb-flow}), $\xi_q\to\xi_1=\xi$ becomes
the classical Reeb field, and the Heisenberg commutator
$[\hat{X},\hat{P}]_q\to 0$
(Remark~\ref{rem:heisenberg-commutator}): the Bring variety
collapses to a non-rotating algebraic surface.

\medskip
\noindent\textbf{Exceptional divisor as mixing interface.}\;
Blowing up a spacetime defect $z_c$ where chiral sheets merge
(Remark~\ref{rem:brings-smoothness}) replaces the singular point with
an exceptional divisor $E_c\cong\mathbb{CP}^1\cong S^2$.  This
sphere acts as the boundary between the macroscopic coordinates
$(x_0,x_1)$ and the internal chiral modes $(w_1,w_2,w_3)$: the
self-intersections of the chiral fibres are resolved, and the
$C_5$ phase rotation now circulates \emph{around} the exceptional
divisor rather than passing through a singularity.

\medskip
\noindent\textbf{Reeb vector as torsion field.}\;
On the blown-up surface, the Reeb vector field~$\xi_q$ is the
infinitesimal generator of the $U(1)$ action that rotates the chiral
fibre above the exceptional divisor.  Concretely, the
$q$-deformed Reeb flow $\phi_t^q$
(Definition~\ref{def:q-reeb-flow}) restricted to $E_c$ generates
the Hopf fibration $S^3\to S^2\cong E_c$, and the substrate
parameter~$q$ controls the pitch of this rotation.  The vector
field~$\xi_q$ thereby provides the ``time-like'' direction within the
chiral space: it is the flow that allows the spectral modes~$w_k$ to
execute the spinning phase rotation
$w_k\mapsto\zeta^k\,w_k$.

\medskip
\noindent\textbf{Global picture.}\;
The Bring surface provides the genus-$4$ topology---four handles
through which the chiral tubes thread.  The Reeb vector field gives
the longitudinal flow along these handles.  The quantum torsion~$q$
modulates the pitch of this flow: when $q\to\infty$ (deep-quantum,
Big Bang) the spinning is maximal and the Heisenberg uncertainty
$[\hat{X},\hat{P}]_q\to J$; when $q\to 1$ (heat death) the spinning
ceases, the Reeb field trivialises, and the topology reduces to a
static cone.
\end{remark}


\begin{remark}[Generation-Dependent Seiberg-Witten Differentials]
The classical Seiberg-Witten integrals over the $A$ and $B$ cycles [cite: 954] fail to distinguish lepton generations. In the QOU framework, the vertical degree of the Reeb decomposition lifts the geometry from the base state bundle $E_1$ to the jet prolongation $J^1E_1$. Consequently, the meromorphic 1-form governing the chiral dynamics is strictly generation-dependent.

For the electron (vertical degree 1), the mass is generated by the baseline Seiberg-Witten differential smoothly integrating over the resolved exceptional divisor $E_c$:
$$a_D^{(e)} = \oint_{E_c} \lambda_{SW}^{(0)} = \oint_{E_c} X dz = \oint_{E_c} (c + z) dz$$

For the muon (vertical degree 2), the gauge field is lifted to the 1-jet bundle, forcing the integration of the momentum gradient. The prolonged differential is $\lambda_{SW}^{(1)} = dX$. Substituting the implicit differentiation of the color tube cubic yields the explicit dual mass integral:
$$a_D^{(\mu)} = \oint_{E_c} \lambda_{SW}^{(1)} = \oint_{E_c} \frac{4z^3+3z^2+2z+1}{c \left( 3c^2 + 2(4z+1)c + (6z^2+3z+1) \right)} dz$$

At the topological defect $z_c$ where chiral sheets merge, the discriminant vanishes and $\partial F / \partial c \to 0$. This induces a severe pole in the prolonged differential $\lambda_{SW}^{(1)}$. By Cauchy's Residue Theorem, the integral over the exceptional divisor captures the residue of this jet-bundle gradient. This geometrically inflates the trapped quantum dilogarithm volume, producing the mass hierarchy $m_\mu \gg m_e$ purely from the topology of the blowup, without requiring arbitrarily high crossing numbers.
\end{remark}

\begin{proposition}[Abelianisation of the MI integral on exceptional divisors]
\label{prop:mi-sw-abelianisation}
Let $E_c \cong \mathbb{CP}^1$ be the exceptional divisor at a branch
point $z_c$ (Remark~\ref{rem:brings-smoothness}), and let
$(\nabla_A q,\; L \wedge F^+)$ be the connection--curvature pair of the
mass identity (Definition~\ref{def:mass-identity}).  Then for each
vertical degree $n \geq 0$ in the Reeb decomposition, the restriction
of the nonabelian multiplicative integral to $E_c$ abelianises to the
$n$-th prolonged Seiberg--Witten period:
\begin{equation}\label{eq:mi-sw-abelianisation}
  \bigl\lVert\,
  \mathrm{MI}\!\!\int_{E_c}
  \bigl(\nabla_A q,\; L \wedge F^+\bigr)
  \,\bigr\rVert_{J^n}
  \;=\;
  \bigl\lvert a_D^{(n)}\bigr\rvert
  \;=\;
  \biggl\lvert\oint_{E_c}\lambda_{\mathrm{SW}}^{(n)}\biggr\rvert,
\end{equation}
where $\lVert\cdot\rVert_{J^n}$ denotes the Killing-form norm projected
onto the $n$-th jet component, and $\lambda_{\mathrm{SW}}^{(n)}$ is the
$n$-th prolonged Seiberg--Witten differential.
\end{proposition}

\begin{proof}
The proof proceeds in three steps.

\emph{Step~1: Abelianisation of the connection.}\;
On the exceptional divisor $E_c \cong \mathbb{CP}^1$ the principal
$\mathrm{SU}(2)$-bundle $P$ trivialises (every $\mathrm{SU}(2)$-bundle
over $S^2$ is classified by $\pi_1(\mathrm{SU}(2)) = 0$).  The
connection $\nabla_A$ therefore reduces to a $\mathfrak{su}(2)$-valued
$1$-form $A|_{E_c}$.  By the weight-space decomposition
(Remark~\ref{rem:weight-space-decomposition}), the restriction splits as
\[
  A\big|_{E_c}
  \;=\;
  A^0 \otimes H
  \;+\;
  A^+ \otimes E_+
  \;+\;
  A^- \otimes E_-,
\]
where $H, E_\pm$ are the standard generators of $\mathfrak{su}(2)$.
The Cartan component $A^0$ is a $\mathrm{U}(1)$-valued $1$-form on
$E_c$.  Since $E_c \cong \mathbb{CP}^1$ the spectral curve
(Remark~\ref{rem:chiral-bundles}) restricts to a single sheet, and the
Higgs-field eigenvalue is the chiral coordinate~$X$.  The abelian
component therefore satisfies $A^0 = X\,dz = \lambda_{\mathrm{SW}}$,
recovering the Seiberg--Witten
differential~\eqref{eq:seiberg-witten-differential}.

\emph{Step~2: Fake-flatness and the $2$-form datum.}\;
The fake-flatness condition $F_A = B$~\eqref{eq:fake-flatness} identifies
the curvature $2$-form $L \wedge F^+$ on $E_c$ with the
$\mathfrak{su}(2)$-valued $2$-form $B|_{E_c}$.  The simplicity constraint
forces $B|_{E_c} = F_A^+|_{E_c}$, which on the $2$-sphere $E_c$ is
proportional to the volume form $\mathrm{vol}_{E_c}$.  Consequently
the pair $(\nabla_A q,\; L \wedge F^+)$ restricted to $E_c$ is
\emph{abelian}: both components lie in the Cartan subalgebra after
projection onto the weight-$0$ sector, and the nonabelian MI integral
reduces to the ordinary (abelian) multiplicative integral, which equals
the exponential of the classical contour integral.  Taking the norm
yields the absolute value of the period.

\emph{Step~3: Jet prolongation and vertical degree.}\;
At vertical degree $n = 0$ the Killing-form projection onto $J^0$
recovers the base differential $\lambda_{\mathrm{SW}}^{(0)} = X\,dz$
and the electron period $a_D^{(e)}$.  At vertical degree $n \geq 1$ the
jet prolongation (Proposition~\ref{prop:jet-prolongation}) lifts the
connection to $J^n E$, replacing $X\,dz$ with the $n$-th prolonged
differential $\lambda_{\mathrm{SW}}^{(n)}$: the $1$-jet component is
$dX$ (the implicit derivative of the color tube
cubic~\eqref{eq:color-tube-cubic}), and higher jets follow by induction.
The MI integral at jet level~$n$ is the abelian integral of
$\lambda_{\mathrm{SW}}^{(n)}$ over $E_c$ by the same argument as
Steps~1--2 applied to the prolonged connection $A'$ on $J^n E$.
\end{proof}

% ── Quantum Brings ───────────────────────────────────────────────────────────

\section{Quantum Brings}
\label{sec:quantum-brings}

% TODO: Develop the quantum version of the commutator relations in the
% group relations / Bring's formulas.  In particular:
% - Replace the classical commutator [e_a,e_b] = epsilon_{abc} e_c with
%   the q-deformed commutator from U_q(su(2)).
% - Describe how the Bring surface equations transform under the
%   q-deformed action.
% - Relate the non-commutative geometry to the quantum substrate
%   parameter q(p).
% - Connect to the universal enveloping algebra framework in
%   Section~\ref{sec:enveloping-algebras}.

% ── Other Euler characteristics ──────────────────────────────────────────────

\section{Other Euler characteristics}
\label{sec:other-euler}

% TODO: Investigate the Bring surface and its analogues for Euler
% characteristics other than chi = -6.
% - Classify surfaces for chi = -2, -4, -8, ...
% - Determine which automorphism groups arise.
% - Connect to the varying Euler characteristic discussion in
%   Section~\ref{sec:varying-euler}.


% ── Section 4: Knots on Surfaces of Genus at Most Four ───────────────────────

\section{Knots on Surfaces of Genus at Most Four}
\label{sec:knots-genus-leq-4}

\subsection{Particle ansatz}
\label{subsec:particle-ansatz}

\begin{remark}[Knot--particle correspondence]
\label{rem:knot-particle-correspondence}
The Bring hierarchy provides the geometric scaffolding for a knot--particle
correspondence: stable particle states of the QOU are identified with
isotopy classes of knots (or links) that can be drawn on a surface of
genus at most $4$ --- i.e.\ knots $K$ with Seifert genus
$g_s(K) \le 4$.  The embedding surface is one of the $B_n$:
\begin{itemize}
  \item $g_s(K)=0$: $K$ embeds on $B_1 = \mathbb{P}^1$; corresponds to
        massless/trivial states (neutrinos).
  \item $g_s(K)=1$: $K$ embeds on the elliptic fiber of $B_3$;
        corresponds to first-generation fermions (electron, up/down
        quark).
  \item $g_s(K)=2$: $K$ embeds on $B_2$; second-generation
        fermions (muon, strange/charm quark).
  \item $g_s(K)=3,4$: $K$ embeds on $B_3$; third-generation
        fermions (tau, bottom/top quark) and heavier states.
\end{itemize}
Composite particles (hadrons, nuclei) correspond to connected sums
$K_1 \mathbin{\#} K_2$ whose Seifert genus equals $g_s(K_1)+g_s(K_2)$.
\end{remark}

\begin{remark}[Knots embeddable in Bring's surface]
\label{rem:knots-embeddable}
A knot $K \subset S^3$ embeds on a Seifert surface of genus $g$ if and
only if $g_s(K) \le g$.  Since Bring's surface $B_3$ has genus $4$, the
knots relevant to the QOU are exactly those with $g_s(K) \le 4$.  Every
torus knot $T(p,q)$ with Seifert genus
$g_s(T(p,q)) = \tfrac{(p-1)(q-1)}{2} \le 4$ embeds.  The fundamental
knots for the QOU particle spectrum are:
\begin{itemize}
  \item $0_1$ (unknot), $g_s=0$: neutrino.
  \item $2^2_1$ (Hopf link), $g_s=0$: gauge boson (photon).
  \item $3_1 = T(2,3)$, $g_s=1$: quark (trefoil, first generation).
  \item $4_1$, $g_s=1$: amphichiral first-generation state.
  \item $5_1 = T(2,5)$, $g_s=2$: second-generation fermion.
  \item $5_2$, $g_s=1$: excited first-generation state.
  \item $T(3,4)$, $g_s=3$: third-generation fermion.
  \item $T(2,7)$, $g_s=3$: alternative third-generation state.
  \item $T(3,5)$, $g_s=4$: maximal-genus fermion on $B_3$.
  \item $T(2,9)$, $g_s=4$: tau-like state at the boundary.
\end{itemize}
Composite knots $3_1 \mathbin{\#} 3_1$ ($g_s=2$) and
$3_1 \mathbin{\#} 4_1$ ($g_s=2$) correspond to two-quark bound
states.
\end{remark}

\subsection{Knot invariant table}
\label{subsec:knot-invariant-table}

The following table lists invariants for each knot.  Columns:
\textbf{Name}: standard Rolfsen/Conway notation.
\textbf{$\Delta(t)$}: Alexander polynomial (symmetric, $\Delta(1)=1$).
\textbf{$\nabla(z)$}: Conway polynomial ($\Delta(t)=\nabla(t^{1/2}-t^{-1/2})$).
\textbf{$V(t)$}: Jones polynomial.
\textbf{$\mathrm{vol}$}: hyperbolic volume of the complement ($0$ for
non-hyperbolic).
\textbf{A/C}: amphichiral~(A) or chiral~(C).
\textbf{$g_s$}: Seifert (minimal) genus.
\textbf{Tw}: twist knot (Y/N).
\textbf{Bl}: algebraic link / singularity blowup (Y/N).

\bigskip
{\renewcommand{\arraystretch}{1.4}
\noindent
\begin{tabular}{llllrcccl}
\hline
\textbf{Name} & \textbf{$\Delta(t)$} & \textbf{$\nabla(z)$} & \textbf{$V(t)$}
  & \textbf{vol} & \textbf{A/C} & \textbf{$g_s$} & \textbf{Tw} & \textbf{Bl} \\
\hline
$0_1$ unknot
  & $1$
  & $1$
  & $1$
  & $0$
  & A & $0$ & Y & Y \\[2pt]
$2^2_1$ Hopf
  & $-1$ (link)
  & $z$
  & $-t^{-1/2}-t^{-5/2}$
  & $0$
  & C & $0$ & --- & Y \\[2pt]
$3_1$ trefoil
  & $t^{-1}-1+t$
  & $1+z^2$
  & $-t^{-4}+t^{-3}+t^{-1}$
  & $0$
  & C & $1$ & Y & Y \\[2pt]
$4_1$ fig.-eight
  & $-t^{-1}+3-t$
  & $1-z^2$
  & $t^2-t+1-t^{-1}+t^{-2}$
  & $2.029$
  & A & $1$ & Y & N \\[2pt]
$5_1 = T(2,5)$
  & $t^{-2}-t^{-1}+1-t+t^2$
  & $1+3z^2+z^4$
  & $-t^{-4}+t^{-3}-t^{-2}+t^{-1}-1+t$
  & $0$
  & C & $2$ & N & Y \\[2pt]
$5_2$
  & $2t^{-1}-3+2t$
  & $1+2z^2$
  & $-t^6+t^5+t^3$
  & $2.828$
  & C & $1$ & Y & N \\[2pt]
$T(3,4)$
  & $t^{-3}-t^{-2}+1-t^2+t^3$
  & $1+2z^2+z^4$
  & see~\eqref{eq:jones-T34}
  & $0$
  & C & $3$ & N & Y \\[2pt]
$T(2,7)$
  & $t^{-3}-t^{-2}+t^{-1}-1+t-t^2+t^3$
  & $1+4z^2+3z^4+z^6$
  & see~\eqref{eq:jones-T27}
  & $0$
  & C & $3$ & N & Y \\[2pt]
$T(3,5)$
  & $t^{-4}-t^{-3}+t^{-1}-1+t-t^3+t^4$
  & $1+3z^2+z^4+z^6$
  & see~\eqref{eq:jones-T35}
  & $0$
  & C & $4$ & N & Y \\[2pt]
$T(2,9)$
  & $\sum_{k=-4}^{4}(-1)^k t^k$
  & $1+5z^2+6z^4+3z^6+z^8$
  & see~\eqref{eq:jones-T29}
  & $0$
  & C & $4$ & N & Y \\[2pt]
$3_1 \mathbin{\#} 3_1$
  & $(t^{-1}-1+t)^2$
  & $(1+z^2)^2$
  & $(-t^{-4}+t^{-3}+t^{-1})^2$
  & $0$
  & C & $2$ & N & N \\[2pt]
$3_1 \mathbin{\#} 3_1^*$
  & $(t^{-1}-1+t)^2$
  & $(1+z^2)^2$
  & $(-t^{-4}+t^{-3}+t^{-1})(-t^4+t^3+t)$
  & $0$
  & A & $2$ & N & N \\[2pt]
$3_1 \mathbin{\#} 4_1$
  & $(t^{-1}-1+t)(-t^{-1}+3-t)$
  & $(1+z^2)(1-z^2)$
  & $(-t^{-4}+t^{-3}+t^{-1})(t^2-t+1-t^{-1}+t^{-2})$
  & $2.029$
  & C & $2$ & N & N \\[2pt]
\hline
\end{tabular}}

\bigskip
\noindent\textbf{Jones polynomials for higher torus knots.}\;
Using the Kauffman bracket / HOMFLY skein relation
$t^{-1}V(L_+) - tV(L_-) = (t^{1/2}-t^{-1/2})V(L_0)$:
\begin{align}
  V(T(3,4)) &= t^{-8}+t^{-6}-t^{-5}+t^{-4}-t^{-3}+t^{-2},
    \label{eq:jones-T34}\\[4pt]
  V(T(2,7)) &= -t^{-10}+t^{-9}-t^{-8}+t^{-7}+t^{-5},
    \label{eq:jones-T27}\\[4pt]
  V(T(3,5)) &= -t^{-16}+t^{-15}-t^{-14}+2t^{-13}-t^{-12}+t^{-11}-t^{-10}+t^{-9},
    \label{eq:jones-T35}\\[4pt]
  V(T(2,9)) &= -t^{-14}+t^{-13}-t^{-12}+t^{-11}+t^{-9}.
    \label{eq:jones-T29}
\end{align}

\subsection{Knot group presentations}
\label{subsec:knot-groups}

\begin{remark}[Groups of the fundamental knots]
\label{rem:knot-groups}
The fundamental group $G_K = \pi_1(S^3 \setminus K)$ of each knot in the
particle ansatz has the following presentation.  For torus knots
$T(p,q)$ the group is the \emph{torus knot group}
$\langle a, b \mid a^p = b^q \rangle$.

\medskip
{\renewcommand{\arraystretch}{1.3}
\begin{tabular}{lll}
\hline
\textbf{Knot} & \textbf{Generators} & \textbf{Relations} \\
\hline
$0_1$ & $\langle a \rangle$ & none \\
$2^2_1$ Hopf & $\langle a, b \rangle$ & $ab = ba$ \\
$3_1$ & $\langle a, b \rangle$ & $a^2 = b^3$ \\
$4_1$ & $\langle x, y \rangle$ & $xyx^{-1}yx = yxy^{-1}xy^{-1}$ \\
$5_1 = T(2,5)$ & $\langle a, b \rangle$ & $a^2 = b^5$ \\
$5_2$ & $\langle x, y \rangle$ & $xy^2x = yx^2y$ \\
$T(3,4)$ & $\langle a, b \rangle$ & $a^3 = b^4$ \\
$T(2,7)$ & $\langle a, b \rangle$ & $a^2 = b^7$ \\
$T(3,5)$ & $\langle a, b \rangle$ & $a^3 = b^5$ \\
$T(2,9)$ & $\langle a, b \rangle$ & $a^2 = b^9$ \\
$3_1\mathbin{\#}3_1$ & $\langle a,b,c,d \rangle$
  & $a^2=b^3,\;c^2=d^3,\;a=c_{\mathrm{merid}}$ \\
$3_1\mathbin{\#}3_1^*$ & $\langle a,b,c,d \rangle$
  & $a^2=b^3,\;c^2=d^{-3},\;a=c_{\mathrm{merid}}$ \\
\hline
\end{tabular}}

\medskip\noindent
The figure-eight group $G_{4_1}$ has the equivalent Fox presentation
$\langle x,t \mid txt^{-1}xt = xtx^{-1}tx^{-1}\rangle$ and is the only
hyperbolic knot group in the table, corresponding to the non-zero
hyperbolic volume.  For torus knot groups, the center is generated by
the element $a^p = b^q$ (the longitude), which maps to the Seifert fiber
in the Seifert fibration of the torus knot complement.
\end{remark}

% ── Section 5: Knot Braiding and Skein Relations ──────────────────────────────

\section{Knot Braiding and Skein Relations}
\label{sec:knot-braiding}

\subsection{HOMFLY-PT skein relation}
\label{subsec:homfly-skein}

\begin{definition}[HOMFLY-PT polynomial]
\label{def:homfly}
The \emph{HOMFLY-PT polynomial} $P(K; \ell, m) \in \mathbb{Z}[\ell^{\pm 1}, m^{\pm 1}]$
of an oriented knot or link $K$ is characterised by the skein relation
\begin{equation}\label{eq:homfly-skein}
  \ell\,P(L_+) + \ell^{-1}\,P(L_-) + m\,P(L_0) = 0,
\end{equation}
together with the normalisation $P(\text{unknot}) = 1$.  Here
$L_+$, $L_-$, $L_0$ are three diagrams identical except near one
crossing, where $L_+$ has a positive crossing, $L_-$ a negative
crossing, and $L_0$ the oriented smoothing.

\medskip\noindent
Specialisations recover classical invariants:
\begin{align}
  &\text{Jones:} &
    V(K)(t) &= P(K;\, t^{-1},\, t^{1/2}-t^{-1/2}),
    \label{eq:jones-specialisation}\\[4pt]
  &\text{Alexander:} &
    \Delta(K)(t) &= P(K;\, i,\, i(t^{1/2}-t^{-1/2})),
    \label{eq:alexander-specialisation}\\[4pt]
  &\text{Conway:} &
    \nabla(K)(z) &= P(K;\, i,\, -iz).
    \label{eq:conway-specialisation}
\end{align}
\end{definition}

\begin{definition}[Kauffman bracket skein relation]
\label{def:kauffman-bracket}
The \emph{Kauffman bracket} $\langle D \rangle \in \mathbb{Z}[A^{\pm 1}]$ of
an unoriented link diagram $D$ satisfies
\begin{equation}\label{eq:kauffman-bracket}
  \langle L_+ \rangle
  = A\,\langle L_{0} \rangle
  + A^{-1}\langle L_{\infty} \rangle,
  \qquad
  \langle D \sqcup \bigcirc \rangle = (-A^2 - A^{-2})\langle D \rangle,
  \qquad
  \langle \bigcirc \rangle = 1,
\end{equation}
where $L_0$ and $L_\infty$ are the two smoothings of a positive crossing $L_+$
(the $A$-smoothing and $A^{-1}$-smoothing respectively).
The Jones polynomial is recovered by
$V(K)(t) = (-1)^{w(D)}A^{-3w(D)}\langle D\rangle\big|_{A=t^{-1/4}}$,
where $w(D)$ is the writhe of a diagram $D$ for $K$.
\end{definition}

\subsection{Skein triples and braid words}
\label{subsec:skein-triples}

\begin{remark}[Skein triples for fundamental knots]
\label{rem:skein-triples}
Each knot in the particle ansatz fits into a skein triple
$(L_+, L_-, L_0)$ showing how to reduce it to simpler knots via the
HOMFLY-PT relation~\eqref{eq:homfly-skein}.  The table below lists the
skein triples and braid word presentations for each knot.  All knots
are oriented; the braid word gives the closure $\hat\beta$ realising the
knot as a closed braid.

\medskip
{\renewcommand{\arraystretch}{1.4}
\noindent\begin{tabular}{llll}
\hline
\textbf{Knot} & \textbf{Braid word} & \textbf{$L_+$} & \textbf{$L_0$} \\
\hline
$0_1$ unknot
  & $1$ (empty)
  & ---
  & --- \\[2pt]
$2^2_1$ Hopf
  & $\sigma_1^2$
  & $2^2_1$
  & $0_1 \sqcup 0_1$ \\[2pt]
$3_1$ trefoil
  & $\sigma_1^3$ ($B_2$-braid)
  & $3_1$
  & $2^2_1$ \\[2pt]
$4_1$ fig.-eight
  & $\sigma_1\sigma_2^{-1}\sigma_1\sigma_2^{-1}$ ($B_3$)
  & $3_1^*$
  & $2^2_1$ \\[2pt]
$5_1 = T(2,5)$
  & $\sigma_1^5$
  & $5_1$
  & $T(2,4)$ \\[2pt]
$5_2$
  & $\sigma_1^2\sigma_2^{-1}\sigma_1^2\sigma_2^{-1}$ ($B_3$)
  & $4_1$
  & $\text{link}$ \\[2pt]
$T(3,4)$
  & $(\sigma_1\sigma_2)^4$ ($B_3$)
  & $T(3,4)$
  & $T(3,3)$ \\[2pt]
$T(2,7)$
  & $\sigma_1^7$
  & $T(2,7)$
  & $T(2,6)$ \\[2pt]
$T(3,5)$
  & $(\sigma_1\sigma_2)^5$ ($B_3$)
  & $T(3,5)$
  & $T(3,4)$ \\[2pt]
$T(2,9)$
  & $\sigma_1^9$
  & $T(2,9)$
  & $T(2,8)$ \\[2pt]
\hline
\end{tabular}}
\end{remark}

\begin{remark}[Explicit skein computations]
\label{rem:skein-computations}
Applying~\eqref{eq:homfly-skein} at the listed crossing for each knot:

\medskip
\noindent\textbf{Trefoil $3_1$ from unknot via Hopf link.}\;
The positive trefoil $3_1^+$ satisfies
\begin{equation}\label{eq:skein-trefoil}
  \ell\,P(3_1^+) + \ell^{-1}\,P(0_1) + m\,P(2^2_1) = 0.
\end{equation}
With $P(0_1)=1$ and $P(2^2_1) = -\ell(m^{-1}+m) - \ell^{-1}m^{-1}$
(two-component unlink skein), one recovers
\[
  P(3_1^+; \ell, m) = -\ell^{-4}+2\ell^{-2}+\ell^{-2}m^2.
\]

\noindent\textbf{Figure-eight $4_1$ as amphichiral fixed point.}\;
The skein relation at the central crossing gives
\begin{equation}\label{eq:skein-fig8}
  \ell\,P(L_+) + \ell^{-1}\,P(L_-) + m\,P(L_0) = 0,
\end{equation}
where $L_+ = 3_1^*$ (negative trefoil) and $L_- = 3_1^+$ (positive trefoil).
Since $P(4_1;\ell,m) = P(4_1;\ell^{-1},m)$ (amphichirality),
the equation constrains $P(4_1)$ to be symmetric in $\ell\leftrightarrow\ell^{-1}$.

\noindent\textbf{Torus knot $T(2,2k+1)$ recursion.}\;
The skein relation on the last positive crossing of $T(2,2k+1)$ gives
\begin{equation}\label{eq:skein-torus-2n}
  \ell\,P(T(2,2k+1)) + \ell^{-1}\,P(T(2,2k-1)) + m\,P(T(2,2k)) = 0,
\end{equation}
with $P(T(2,1))=P(0_1)=1$ and $P(T(2,2))=P(2^2_1)$.
This determines $P(T(2,2k+1))$ inductively for all $k \ge 1$.

\noindent\textbf{Torus knot $T(3,n)$ recursion.}\;
Writing $T(3,n)$ as the closure of $(\sigma_1\sigma_2)^n$ in $B_3$, the skein relation gives
\begin{equation}\label{eq:skein-torus-3n}
  \ell\,P(T(3,n)) + \ell^{-1}\,P(T(3,n-2)) + m\,P(L_n) = 0,
\end{equation}
where $L_n$ is an intermediate link whose HOMFLY polynomial is computed
from lower-crossing diagrams.
\end{remark}

\begin{remark}[Jones skein relation for each knot]
\label{rem:jones-skein-per-knot}
Specialising $\ell = t^{-1}$, $m = t^{1/2}-t^{-1/2}$ in
\eqref{eq:homfly-skein} gives
$t^{-1}V(L_+) - tV(L_-) = (t^{1/2}-t^{-1/2})V(L_0)$.
The table below shows the skein relation at one distinguished crossing
for each knot, together with the resulting Jones computation:

\medskip
{\renewcommand{\arraystretch}{1.4}
\noindent\begin{tabular}{lll}
\hline
\textbf{Knot $K$} & \textbf{Jones skein identity} & \textbf{Result} \\
\hline
$0_1$
  & $V(0_1) = 1$ (base case)
  & $1$ \\[2pt]
$2^2_1$
  & $t^{-1}V(2^2_1) - t\,V(0_1\sqcup 0_1) = (t^{1/2}-t^{-1/2})V(0_1)$
  & $-t^{-1/2}-t^{-5/2}$ \\[2pt]
$3_1$
  & $t^{-1}V(3_1) - t\,V(0_1) = (t^{1/2}-t^{-1/2})V(2^2_1)$
  & $-t^{-4}+t^{-3}+t^{-1}$ \\[2pt]
$4_1$
  & $t^{-1}V(3_1^*) - t\,V(3_1) = (t^{1/2}-t^{-1/2})V(2^2_1)$
  & $t^2-t+1-t^{-1}+t^{-2}$ \\[2pt]
$5_1$
  & $t^{-1}V(5_1)-t\,V(3_1)=(t^{1/2}-t^{-1/2})V(T(2,4))$
  & see~\eqref{eq:jones-T35} \\[2pt]
$5_2$
  & $t^{-1}V(4_1)-t\,V(5_2^*)=(t^{1/2}-t^{-1/2})V(L)$
  & $-t^6+t^5+t^3$ \\[2pt]
$T(3,4)$
  & $t^{-1}V(T(3,4))-t\,V(T(3,2))=(t^{1/2}-t^{-1/2})V(T(3,3))$
  & see~\eqref{eq:jones-T34} \\[2pt]
$T(2,7)$
  & $t^{-1}V(T(2,7))-t\,V(T(2,5))=(t^{1/2}-t^{-1/2})V(T(2,6))$
  & see~\eqref{eq:jones-T27} \\[2pt]
$T(3,5)$
  & $t^{-1}V(T(3,5))-t\,V(T(3,3))=(t^{1/2}-t^{-1/2})V(T(3,4))$
  & see~\eqref{eq:jones-T35} \\[2pt]
$T(2,9)$
  & $t^{-1}V(T(2,9))-t\,V(T(2,7))=(t^{1/2}-t^{-1/2})V(T(2,8))$
  & see~\eqref{eq:jones-T29} \\[2pt]
\hline
\end{tabular}}
\end{remark}

\begin{remark}[$q$-deformed skein relation in the QOU]
\label{rem:q-deformed-skein}
In the QOU framework, the standard Jones skein relation is replaced by
a $q$-deformed version incorporating the substrate parameter
$q(p)$ (Definition~\ref{def:substrate-parameter}):
\begin{equation}\label{eq:q-skein}
  q^{-1}\,V_q(L_+) \;-\; q\,V_q(L_-) \;=\;
  (q^{1/2}-q^{-1/2})\,V_q(L_0),
\end{equation}
where $V_q(K)$ is the $q$-deformed Jones polynomial evaluated at the
substrate parameter $q(p)$ of the QOU point $p \in B_3$.  At
$q\to 1$ (classical limit, heat death) equation~\eqref{eq:q-skein}
reduces to the ordinary Jones skein relation.  At $q\to\infty$
(deep-quantum, Big Bang) the $q$-deformed $R$-matrix of
$\mathcal{U}_q(\mathfrak{su}(2))$
(Remark~\ref{rem:r-matrix}) furnishes the twisted exchange operator,
and the braiding becomes maximal.

\medskip\noindent
The $q$-deformation parameter couples the knot invariant to the
topology of the Bring hierarchy: the substrate parameter $q(p)$ is
maximal at the six branch points (where $B_3$ and $B_2$ have their
most complex fiber structure), and the $q$-deformed skein relation
therefore encodes the transition between hierarchy levels as a
deformation of the classical braid relations.
\end{remark}
